\documentclass[oneside]{amsart}
\usepackage[utf8]{inputenc}
\usepackage[T1]{fontenc}
\usepackage{amsmath,amssymb,amsthm,bbm,float,graphicx,geometry,lmodern,mathtools,parskip,setspace,subcaption,diagbox}

\usepackage{authblk}

\usepackage{mathtools}
\mathtoolsset{showonlyrefs}

\usepackage[colorlinks]{hyperref}
\usepackage{mathrsfs}
\usepackage{enumerate}
\usepackage{nicefrac}
\usepackage{todonotes}
\usepackage{comment}
\usepackage[notref,notcite]{showkeys}
\newcommand{\e}{\mathrm e}

%%%Tikz
\usepackage{tikz}
\usetikzlibrary{backgrounds}
\usetikzlibrary{patterns}
\usetikzlibrary{positioning, shapes.geometric}

%%% Captionsize
\usepackage{caption}
%\usepackage{subcaption}
\captionsetup{font=footnotesize}
\usepackage{graphicx}
%%%%%%Formatting%%%%%%
\graphicspath{{Images/}}%put any images in this file
\parindent0pt%paragraph indent
\allowdisplaybreaks%breaks long equations onto multiple pages is needed
%%%%%%Commands%%%%%%
\DeclarePairedDelimiter{\set}{\lbrace}{\rbrace}
\newcommand{\N}{\mathbb{N}}
\DeclareMathOperator{\argmax}{arg \ max}
%%%%%%Environments%%%%%%

\newtheorem{theorem}{Theorem}[section]
\newtheorem{lemma}[theorem]{Lemma}
\newtheorem{corollary}[theorem]{Corollary}
\newtheorem{prop}[theorem]{Proposition}
\newtheorem{conj}[theorem]{Conjecture}

\theoremstyle{definition}
\newtheorem{definition}[theorem]{Definition}
\newtheorem{remark}[theorem]{Remark}

%%%%%%%% symbols and short cuts%%%%%%%
\renewcommand{\P}{\mathbb{P}}
\newcommand{\x}{\mathbf{x}}
\newcommand{\y}{\mathbf{y}}
\newcommand{\1}{\mathbbm{1}}
\newcommand{\G}{\mathscr{G}}
%\renewcommand{\H}{\mathscr{H}}
\newcommand{\C}{\mathscr{C}}
\newcommand{\X}{\mathbf{X}}
\newcommand{\E}{\mathbb{E}}
\newcommand{\0}{\mathbf{o}}
\newcommand{\R}{\mathbb{R}}
\renewcommand{\d}{\mathrm{d}}
\newcommand{\z}{\mathbf{z}}
\newcommand{\w}{\mathbf{w}}
\renewcommand{\v}{\mathbf{v}}
\renewcommand{\u}{\mathbf{u}}
\newcommand{\Z}{\mathbb{Z}}
\newcommand{\T}{\mathcal{T}}
\newcommand{\U}{\mathcal{U}}
\newcommand{\cX}{\mathcal{X}}
\newcommand{\cG}{\mathcal{G}}


%Commands Bas
\newcommand*{\be}{\begin{equation}}
\newcommand*{\ee}{\end{equation}}
\newcommand*{\ba}{\begin{aligned}}
\newcommand*{\ea}{\end{aligned}}

\newcommand{\D}{\mathbb D}
\newcommand{\cB}{\mathcal B}
\newcommand{\cP}{\mathcal P}
\newcommand{\cU}{\mathcal U}
%\newcommand{\ind}{\mathbbm 1}
\newcommand{\wt}{\widetilde}
\newcommand{\cR}{\mathcal R}
\newcommand{\sm}{\mathrm{sm}}
\newcommand{\basc}[1]{{\color{blue}{ \bf [~Bas:\ }\emph{#1}\textbf{~]}}}

%Packages Bas
%\usepackage{enumitem}
\usepackage{bbm}

\numberwithin{equation}{section}

%%%%%%Bibliography%%%%%%
\usepackage[style = numeric, sorting=nyt, url = false, abbreviate=false, maxbibnames=9, sortcites=true, doi = true, backend = biber, giveninits = true]{biblatex}
\renewbibmacro{in:}{\ifentrytype{article}{}{\printtext{\bibstring{in}\intitlepunct}}}
\bibliography{bib.bib}

\title{Percolation in lattice $k$-neighbor graphs}
\author[1,2]{Benedikt Jahnel} \author[2]{Jonas K\"oppl}
\author[3,4]{Bas Lodewijks} \author[5,6]{Andr\'as T\'obi\'as}
\affil[1]{TU Braunschweig, Germany}
\affil[2]{WIAS, Berlin, Germany}
\affil[3]{Univ.\ Jean Monnet, Saint-Etienne, France}
\affil[4]{Institut Camille Jordan, Lyon and Saint-Etienne, France}
\affil[5]{Alfr\'ed R\'enyi Institute of Mathematics, Budapest, Hungary}
\affil[6]{Budapest University of Technology and Economics, Hungary}
\date{\today}

\begin{document}

\maketitle

\begin{abstract}
%1:
We define a random graph obtained by connecting each point of $\Z^d$ independently to a fixed number $1 \leq k \leq 2d$ of its nearest neighbors via a directed edge. 
%2:
We call this graph the \emph{directed $k$-neighbor graph}. Two natural associated undirected graphs are the \emph{undirected} and the \emph{bidirectional} $k$-neighbor graph, where we connect two vertices by an undirected edge whenever there is a directed edge in the directed $k$-neighbor graph between the vertices in at least one, respectively precisely two, directions. 
%3: 
For these graphs we study the question of percolation, i.e., the existence of an infinite self-avoiding path. 
%4:
Using different kinds of proof techniques for different classes of cases, we show that for $k=1$ even the undirected $k$-neighbor graph never percolates, while the directed $k$-neighbor graph percolates whenever $k \geq d+1$, $k \geq 3$ and $d \geq 5$, or $k \geq 4$ and $d=4$.
%5:
We also show that the undirected $2$-neighbor graph percolates for $d=2$, the undirected $3$-neighbor graph percolates for $d=3$, and we provide some positive and negative percolation results regarding the bidirectional graph as well.
%6:
A heuristic argument for high dimensions indicates that this class of models is a natural discrete analogue of the $k$-nearest-neighbor graphs studied in continuum percolation, and our results support this interpretation.
\end{abstract}

\emph{Keywords and phrases.} Lattice $k$-neighbor graphs, directed $k$-neighbor graph, undirected $k$-neighbor graph, bidirectional $k$-neighbor graph, 1-dependent percolation, oriented percolation, negatively correlated percolation models, connective constant, planar duality, coexistence of phases.

\smallskip

\emph{MSC 2020.} 60K35, 82B43


\section{Introduction}
The study of $k$-nearest-neighbor type models in continuum percolation was initiated by Häggström and Meester \cite{HaggMee96}, who investigated the so-called \emph{undirected} $k$-nearest-neighbor graph. There, the vertex set of the random graph under consideration is a homogeneous Poisson point process in $\R^d$, and one connects two points by an edge if at least one of them is among the $k$ nearest neighbors of the other. One calls a maximal connected component in this graph a \emph{cluster}, and one says that the graph \emph{percolates} if it exhibits an infinite (or equivalently, unbounded) cluster. Note that the probability of percolation is in $\{0,1\}$ thanks to ergodicity, and whether the graph percolates does not depend on the spatial intensity of the Poisson point process. Hence, the only parameters of the model are $k$ and $d$. It was shown in~\cite{HaggMee96} that this model does not percolate for any dimension if $k = 1$, but it percolates for $k=2$ if $d$ is large enough. Moreover, they proved that for any $d \geq 2$, there exists $k$ such that the graph percolates.

Balister and Bollobás \cite{BalisterBollobas13} investigated three possible senses of percolation in the model where each vertex is connected  to its $k$ nearest neighbors by a \emph{directed} edge: \emph{in-percolation} (resp.~\emph{out-percolation}), which occurs by definition whenever some point of the point process exhibits an infinite incoming (resp.\ outgoing) path ending (resp.\ starting) at it, and \emph{strong percolation}, which means that there exists a strongly connected component in the graph, i.e., a component where from any point there exists a directed path to any other point. Additionally, they introduced another undirected graph, called the \emph{bidirectional} $k$-nearest-neighbor graph, where one connects two vertices whenever they are mutually among the $k$ nearest neighbors of each other. For the two-dimensional case, they verified percolation in the undirected graph for $k \geq 11$, in the directed graph in all the three senses (in-, out- and strong percolation) for $k \geq 13$ and in the bidirectional graph for $k \geq 15$. It is easy to see that these graphs do not percolate in one dimension for any $k$. We observe that it follows from the proof of \cite[Theorem 2.2]{HaggMee96} that for any $d\geq 2$, for $k \geq 2$ sufficiently large, also the bidirectional graph percolates (even though the notion of bidirectional $k$-nearest-neighbor graph does not appear in that paper).

%These $k$-nearest-neighbor percolation models can be seen as an alternative of modeling connectivity in a telecommunication network, as opposed to the classical Gilbert graph, where one connects two points of the Poisson point process by an edge whenever their Euclidean distance is less than a given radius $r>0$. The latter graph was introduced by Gilbert (1961), who showed that for $d =2$, there is a phase transition in the spatial intensity $\lambda$ of the Poisson point process: There exists a critical intensity $\lambda_{\rm c}>0$ such that for $\lambda<\lambda_{\rm c}$, the graph almost surely consists of finite components, while for $\lambda>\lambda_{\rm c}$, the graph \emph{percolates}, i.e., it contains an infinite (equivalently, unbounded) connected component. While the Gilbert graph has unbounded degrees but bounded edge lengths, $k$-nearest-neighbor type models have bounded degrees but arbitrarily long edges. A more realistic model for telecommunications, which is a proper subgraph of both the bidirectional $k$-nearest-neighbor graph for some $k$ and the Gilbert graph for some $r$ is the so-called signal-to-interference plus noise ratio (SINR) graph, see \color{red} references \color{black}.

It was further shown in \cite{BalisterBollobas13} that for $d=2$, with high confidence, the undirected $k$-nearest-neighbor graph percolates if and only if $k \geq 3$, in-, out- and strong percolation in the directed graph occurs if and only if $k\geq 4$, and the bidirectional graph percolates if and only if $k\geq 5$. By ``with high confidence'' they mean the following: They show that if a certain deterministic integral exceeds a certain deterministic value, then their aforementioned assertions are true. However, they were only able to evaluate these integrals via Monte--Carlo methods, which yield an extremely small but positive error probability. Recently, Jahnel and T\'obi\'as \cite{JT22} showed that in the bidirectional graph there is no percolation for $k=2$ in any dimension, even if the underlying point process is not a Poisson point process but an arbitrary deletion-tolerant and stationary point process (satisfying some basic nondegeneracy conditions). Their proof exploits the simplicity of the structure of the bidirectional graph for $k=2$, which has degrees bounded by 2. Proving absence of percolation for $k=3,4$ seems to be out of reach at the moment even in the Poisson case, and it is also not entirely clear whether these assertions hold in all dimensions or only for $d=2$.

To understand the structure and elementary properties of $k$-nearest-neighbor graphs better, in the current paper we take a step back and introduce a discrete counterpart of the (undirected, directed or bidirectional) $k$-nearest-neighbor model. Here, $1 \leq k \leq 2d$ (where $d$ still denotes the dimension), and the directed  $k$-neighbor graph ($k$-DnG) is defined via connecting each vertex of $\Z^d$ to precisely $k$ of its $2d$ nearest neighbors (with respect to the $\ell_1$-distance), which are chosen independently and uniformly for each vertex. If we now put an undirected edge between all nearest-neighbor pairs connected by at least one edge (respectively by two directed edges) of the $k$-DnG, we obtain the undirected (respectively bidirectional) $k$-neighbor graph, abbreviated as $k$-UnG (respectively $k$-BnG). In the $k$-DnG we study the analogue of out-percolation, which is often the most convenient choice, but in- or strong percolation could also be investigated similarly, cf.\ Section~\ref{sec-instrong} for a related discussion. Percolation in the $k$-DnG in our sense also implies percolation in the $k$-UnG, while percolation in the $k$-BnG graph implies percolation in the $k$-DnG. In Section~\ref{sec:motivation} we provide more details regarding the relation between the discrete and continuum nearest-neighbor percolation, which gives motivation for studying this particular model.

Note also that the $k$-BnG is a 1-dependent Bernoulli bond percolation model where each edge is ``open'' (i.e., included in the $k$-BnG) with probability $p=k^2/(4d^2)$. That is, any fixed edge is open with probability $p$, and the events that some edges are open are independent whenever the edges are pairwise non-adjacent. The same holds for the $k$-UnG with $p=(1-k/(4d))k/d$. Let us also mention that these lattice $k$-neighbor graphs present basic examples of {\em negatively correlated} random fields, in contrast to many of the classical models like the random cluster model, see \cite{grimmett2006random}. Indeed, note that the presence of an edge makes the presence of neighboring edges less likely. 

Our main results consider the presence and absence of percolation for all the three variants of the model. The ones corresponding to the $k$-DnG and $k$-UnG are summarized in Table~\ref{table-UD}, colored according to the kinds of techniques used in their proofs. For the $k$-DnG, we show percolation for $k\geq 3$ in $d\geq 5$ dimensions and for $k \geq 4$ also in 4 dimensions (cf.~Theorem~\ref{thm-CoxDurrett}) based on arguments by~\cite{CoxDurrett} on oriented percolation in at least 4 dimensions, combined with some explicit computations in a few low-dimensional cases where the estimates of~\cite{CoxDurrett} are not strong enough yet. We also present an argument showing that for $k \geq d+1$ (also in lower dimensions), percolation holds in a very strong sense: There is an infinite outgoing path based from any fixed vertex with probability 1, see Lemma~\ref{lem:sureperc}.
Almost all of our positive percolation results about the $k$-UnG follow from the same assertion for the corresponding $k$-DnG, the only difference is the case $k=d=2$, where using a Peierls argument involving a dual lattice in the spirit of~\cite{Durr88} and~\cite{JanSwart}, we show percolation for the $k$-UnG (but not for the $k$-DnG), cf.\ Theorem~\ref{theorem-U22}. That theorem also shows that the same holds for $k=d=3$: In this case, we verify percolation restricted to a two-dimensional plane and use the same techniques as for $k=d=2$. For $k=1$, however, even the $k$-UnG contains no infinite connected component. This assertion is Proposition~\ref{prop-U1}, whose proof involves a mass-transport type argument that is somewhat similar in spirit to the analogous proof in the continuum case~\cite{HaggMee96}.

As for the $k$-BnG, we show that it never percolates if $k(k-1) < 2d (2d-1)/c(d)$, where $c(d)$ is the connective constant of $\Z^d$ (whose definition we recall in the beginning of Section~\ref{sec:U}), based on negative correlations in this graph. %This implies that if the $k$-BnG percolates, then $k=\Omega(\sqrt d)$ as $d\to\infty$. 
In particular, for fixed $k$ there is no percolation if $d$ is sufficiently large. Such a phenomenon has not been observed in the continuum setting so far. Another consequence is that for $k=2$ we have no percolation in the $k$-BnG for any $d\geq 2$, which is the analogue of the aforementioned assertion by \cite{JT22} in the continuum case. We also obtain some positive percolation results for the $k$-BnG. First, again via verifying percolation restricted to a two-dimensional plane, we show that the $k$-BnG percolates whenever $k>\sqrt{8/3}d \approx 1.6330d$.
Finally, using recent results by~\cite{improvedLB} on 1-dependent Bernoulli bond percolation models, we deduce that the $k$-BnG percolates for $d$ sufficiently large if $k > 2\sqrt{0.5847}d \approx 1.5293 d$. 
\begin{table}
\centering
\begin{tabular}{|l|c|c|c|c|c|c|c|c|} 
 \hline
\diagbox{$d$}{$k$} & 1 & 2 & 3 & 4 & 5 & 6 & $\geq 7$ \\ \hline 
1 & no/no & yes/yes & - & - & - & - & - \\ \hline
2 & \color{blue}no/no\color{blue} & \color{black}open\color{black}/\color{pink}yes\color{black} & \color{green}yes/yes\color{black} & yes/yes & - & - & - \\ \hline
3 & \color{blue}no/no\color{blue} & \color{black}open/open & \color{black}open/\color{pink}yes & \color{green}yes/yes\color{black} & \color{green}yes/yes\color{black} & yes/yes & - \\ \hline
4 & \color{blue}no/no\color{blue} & \color{black}open/open& \color{black}open/open& \color{red}yes/yes\color{black} & \color{green}y\color{red}es/\color{black}\color{green}y\color{red}es\color{black} & \color{green}y\color{red}es/\color{black}\color{green}y\color{red}es\color{black} & \color{green}y\color{red}es/\color{black}\color{green}y\color{red}es\color{black} \\ \hline
$\geq 5$ & \color{blue}no/no\color{blue} & \color{black}open/open& \color{red}yes/yes & \color{red}yes/yes & \color{red}yes/yes & \color{red}yes/yes & \color{red}yes/yes \\  \hline
\end{tabular}

\caption{Percolation results and open cases for the $k$-DnG/$k$-UnG. In the table, for given $k,d$, `yes' means that the given graph percolates, `no' means that it does not, while `open' means that the given case is open. The coloring subdivides the cases according to the assertions from which they follow. Black: trivial or open, blue: Proposition~\ref{prop-U1} (for the directed case, Lemma~\ref{lem:surenoperc} also suffices)\color{black}, pink: Theorem~\ref{theorem-U22}, green: Lemma~\ref{lem:sureperc},  red: Theorem~\ref{thm-CoxDurrett} (and green-and-red: both of the latter two).}\label{table-UD}
\end{table}

\begin{figure}[ht]
    \centering
    \includegraphics[width=1\textwidth]{Images/Comp_origin_2DnG}
    \caption{Connected component of the origin for the $2$-DnG and $d=2$, restricted to $[-100,100]^2$. The origin is in green, the vertices in the origin's connected component are in red, and those not in the connected component are in blue.}
    \label{fig:2dng}
\end{figure}

\begin{figure}[ht]
    \centering
    \includegraphics[width=1\textwidth]{Images/Comp_origin_3BnG}
    \caption{Connected component of the origin for the $3$-BnG and $d=2$, restricted to $[-100,100]^2$. The origin is in green, the vertices in the origin's connected component are in red, and those not in the connected component are in blue.}
    \label{fig:3bng}
\end{figure}



The similarity between some of our main results and assertions or conjectures for the continuum $k$-nearest-neighbor graphs indicate that, despite the lack of the unbounded-range dependencies that are characteristic for the continuum models, the models studied here are indeed their natural discrete analogue. Even without being motivated by the continuum case, our graphs show interesting phase transitions and have many connections to the existing literature (as listed above), while their definition is very simple, which gives a reason to study them in their own right. Nevertheless, to our best knowledge, these models have not been introduced before, which we found somewhat surprising. Similar discrete models such as corner percolation \cite{cornerperc} and constrained-degree percolation \cite{constraineddegreeperc} are known from the literature, but they all seem to be significantly different from those studied here, and so far they have not proven to be useful to verify (absence of) percolation in the models we consider.

The present paper leaves some cases open, which would be interesting to resolve in the future. As already mentioned, we do not know whether the 2-DnG percolates in any dimension, but numerical simulations suggest that this actually seems to be the case for $d=2$ (see Figure~\ref{fig:2dng}). Such an assertion would be particularly interesting because the nearest-neighbor edges of $\Z^d$ missing from the $d$-DnG also form a $d$-DnG, i.e., if the $d$-DnG percolates, we have a coexistence of infinite clusters of chosen and non-chosen edges, which is generally considered as untypical for a two-dimensional percolation model. We also have two ``sporadic'' cases where the question of percolation in the directed graph is still open: $k=d=3$ and $k=3,d=4$, in the latter case it is also open for the undirected graph.

In Theorem~\ref{theorem-3322} below we show that for all $k, d \geq 1$ we have $\theta^{\rm D} (k+1, d+1) \geq \theta^{\rm D}(k,d)$. This result is purely motivated by the open cases. Namely, it implies that if one can show percolation for the $2$-DnG in 2 dimensions, then the same follows for the $3$-DnG in 3 dimensions, whereas if one can show absence of percolation in the latter graph, it implies absence of percolation in the former one. The relation between percolation in the $2$-DnG for $d=3$ and the $3$-DnG for $d=4$ is analogous. 

We see that for any fixed $k$, there is percolation in the $(2d-k)$-BnG and absence of percolation in the $k$-BnG if $d$ is large enough, but finding the precise order of the smallest $k$ for which percolation holds as $d \to \infty$ is not known yet. According to our results, it is $\Omega(\sqrt d)$ and bounded by $(1.5293+o(1))d$. Also, it would be interesting to know if the $k$-BnG ever percolates for $k \leq d$. Moreover, we do not know if the $3$-BnG percolates in two dimensions. One would (at least naively) imagine that this is the case, and numerical simulations indeed indicate this (see Figure~\ref{fig:3bng}). What also supports this conjecture is that the nearest-neighbor edges of $\Z^2$ missing from the $3$-BnG form a $1$-UnG (cf.~Section~\ref{sec:phases} for more details, where some further open questions are explained from a different perspective), for which we have showed that percolation does not occur. 

The remainder of this paper is organized as follows. In Section~\ref{sec:set} we provide a formal definition of our $k$-DnG,  $k$-UnG, and $k$-BnG models, and we present our results corresponding to these graphs in Sections~\ref{sec:D}, \ref{sec:U}, and \ref{sec:B} along with their proofs, respectively. In Section~\ref{sec:discussion} we discuss our results and mention some related conjectures and open questions. The proof of Lemma~\ref{lemma-largedmon}, a technical result used in the proof of Theorem~\ref{thm-CoxDurrett}, is carried out in the Appendix.

\section{Setting and main results}\label{sec:set}

Consider the $d$-dimensional hypercubic lattice $\mathbb Z^d$. We are interested in the percolation behavior of the {\em $k$-neighbor graph} ($k$-nG) in which each vertex independently chooses uniformly $k\le 2d$ of its $2d$ nearest neighbors. We distinguish three different types of $k$-nGs:
\begin{enumerate}
\item The {\em directed $k$-nG} ($k$-DnG) in which we open a directed edge from the vertex to each of the $k$ chosen neighbors. 
\item The {\em undirected $k$-nG} ($k$-UnG) in which we open an undirected edge from the vertex to each of the $k$ chosen neighbors.
\item The {\em bidirectional $k$-nG} ($k$-BnG) in which we open an undirected edge between two vertices whenever both choose the other one as one of their $k$ neighbors.
%\item The {\em exclusively unidirectional $k$-nG} or ($k$-XnG) in which we open an undirected edge between two vertices whenever exactly one of them chooses the other one as one of its $k$ chosen neighbors. The letter X refers to the ``XOR'' (exlusive ``or'') in the edge-drawing rule.
\end{enumerate}

We are interested in the behavior of the percolation probabilities 
\begin{equation}\label{thetaD}
\theta^{\rm D}(k,d)=\mathbb P(o\rightsquigarrow\infty \text{ in the $k$-DnG on $\mathbb Z^d$}),
\end{equation}
respectively $\theta^{\rm U}(k,d)$ and $\theta^{\rm B}(k,d)$, %and $\theta^{\rm X}(k,d)$, 
where $o\rightsquigarrow\infty$ represents the event that there exists a path along open edges from the origin to infinity. Let us note that the percolation probabilities are clearly non-decreasing functions in $k$. 

Let us mention that the outdegree in $k$-DnG is $k$ almost surely. However, the degree in $k$-BnG is a binomial random variable with parameters $k$ and $k/2d$, i.e., the expected degree is given by $k^2/2d$ which is less than $k$, unless $k=2d$. In case of the $k$-UnG, the degree is distributed according to $k+D$, where $D$ is a binomial random variable with parameters $2d-k$ and $k/2d$, and hence, the expected degree is given by $k(4d-k)/(2d)$.

\subsection{Results for the directed $k$-neighbor graph}\label{sec:D}
Let us start with our results for the directed case. First, choosing only one neighbor is never enough for percolation. 
\begin{lemma}\label{lem:surenoperc}
We have for all $d\ge 1$ that $\theta^{\rm D}(1,d)=0$.
\end{lemma}
\begin{proof}
Indeed, if there exists a path starting from the origin and reaching an $\ell_1$-distance $n$ from the origin, then this path is unique and has at least $n$ steps. But such a path exists with probability at most $(2d)^{-n}$, which converges to zero as $n$ tends to infinity.
\end{proof}
%Note that the $k$-nnG has a natural dual in which each existing directed edge is removed and all missing edges are present. This graph is edge-disjoint but not vertex disjoint and it is in distribution equivalent to the $(2d-k)$-nnG. 
% \begin{lemma}\label{lemma2}
% We have that $\theta(2d-1,d)>0$.
% \end{lemma}
% \begin{proof}
% Note that by ergodicity, since $\theta(1,d)=0$ in the $(2d-1)$-nnG in $d$ dimensions, there exist only finite components in the $1$-nnG in $d$ dimensions. But this means that the $(2d-1)$-nnG in $d$ dimensions percolates and thus $\theta(2d-1,d)>0$. 
% \end{proof}
On the other hand, sufficiently large $k$ leads to an almost-sure existence of a path connecting the origin to infinity. 
\begin{lemma}\label{lem:sureperc}
    We have that $\theta^{\rm D}(k,d)=1$ for all $k\ge d+1$.
\end{lemma}
In particular, from Lemma~\ref{lem:surenoperc} and Lemma~\ref{lem:sureperc} we see that for $d=2$ the critical $k$ is given by $2$. Simulations suggest that $\theta^{\rm D}(2,2)>0$, see, e.g., Figure~\ref{fig:2dng}.
\begin{proof}
    We will use a growth argument and show that the maximal distance to the origin is strictly increasing between generations. To make this precise, denote by $G_n$ the new vertices discovered in the $n$-th step, where we start with $G_0 = \{0\}$ and at every step $n \to n+1$, each vertex in $G_n$ chooses $k$ of its neighbors uniformly at random as successors and $G_{n+1}$ is then the set of all potential successors that have not been discovered before for any $m  \leq n$. For $x\in\Z^d$, we let $\|x\|_1$ denote the $\ell_1$-distance between $x$ and the origin. We show by induction that 
    \begin{align}\label{eqn:growth-property}
         \forall n  \in \N:   \quad \max_{x \in G_n}\|x\|_1 < \max_{x \in G_{n+1}}\|x\|_1.  
    \end{align}
    This then implies that there exists an infinitely long directed path starting at the origin. 
    
    Indeed, for $n=0$ the inequality in~\eqref{eqn:growth-property} is clear. For the induction step, let $x \in G_n$ be a vertex (possibly non unique) that achieves the maximum $\ell_1$ distance. Then $x$ has $d$ neighbors $y$ with $\|x\|_1 < \|y\|_1$. If $k  \geq 2d-d+1=d+1$, then we necessarily have to choose one of these $y$ as a potential vertex for the new generation. This $y$ cannot have been in any of the previous generations by the induction hypothesis. Therefore we have $y \in G_{n+1}$ and~\eqref{eqn:growth-property} follows. 
\end{proof}

Let us turn our attention towards the intermediate supercritical percolation phase. 
% For this, consider two independent random walks $S=(S_n)_{n\in\N_0}$ and $S'=(S'_n)_{n\in\N_0}$ on $\Z^d$ with $S_0=e_1$ and $S'_0=e_2$, where $(e_1,\dots, e_d)$ denote the unit vectors. The walks move as simple random walks in the positive direction (the first quadrant), i.e., $S_{n+1}=S_n+e_i$ where $i$ is uniformly chosen from $\{1,\dots ,d\}$. We define by
% \begin{align*}
% \rho(d)=\P(\exists m\ge 1\colon S_m=S'_m)\in(0,1]
% \end{align*}
% the meeting probability of the walkers. 
We have the following result.
\begin{theorem}\label{thm-CoxDurrett}
$\theta^{\rm D}(k,d)>0$ holds if $d \geq 4$ and $k \geq 4$, or if $d \geq 5$ and $k=3$. 
\end{theorem}
Clearly, using Lemma~\ref{lem:sureperc}, for $k \geq 4$ we even have percolation for any $d \geq 2$, but we did not include this in the statement of the theorem in order to separate assertions with different proof techniques.
\begin{proof}[Proof of Theorem~\ref{thm-CoxDurrett}]
 We follow \cite[Section 2]{CoxDurrett}. Let $\mathcal R_n$ denote the set of directed paths from the origin to level $n$ (that is, all vertices at $\ell_1$ distance $n$) with strictly increasing $\ell_1$ distance in the first quadrant, and let $N_n$ be the number of open paths in $\mathcal R_n$. Then $W_n=(2/k)^nN_n$ is a martingale with respect to $\mathcal F_n$, the sigma-algebra generated by the first $n$ steps of a simple random walk $T$ with $T_0=o$ and $T_{n+1}=T_n+e_i$ where $i$ is uniformly chosen from $\{1,\dots ,d\}$. Indeed, using the fact that the expected number of outgoing (i.e., w.r.t.\ $\ell_1$ distance increasing) edges  is given by the expectation of a hypergeometric random variable, i.e., 
\begin{align*}
\sum_{\ell=0}^k \ell\binom{d}{\ell}\binom{d}{k-\ell}\Big/\binom{2d}{k}=\frac{k}{2},
\end{align*}

\begin{comment} 
\basc{Observe that this is only true for vertices $v$ such that the number of neighbors with larger (resp.\ smaller) $\ell_1$-distance equals $d$ (resp.\ $d$). There are, however, vertices for which this is not the case. In any dimension $d\geq 2$ and any $n\in\N$, there are always $2d$ vertices at distance $n$ such that they have $2d-1>d$ many neighbors at $\ell_1$-distance $n+1$ and only $1$ neighbor at $\ell_1$-distance $n-1$. In particular, when $d=k=2$,
\be 
N_1=k \text{ almost surely, and } \E[N_2]=12\Big(\frac{k}{2d}\Big)^2=3\neq 2. 
\ee
Similarly, 
\be \ba
\E[N_3]&=4\cdot 3\Big(\frac{k}{2d}\Big)^3+4\cdot 2\cdot 2\Big(\frac{k}{2d}\Big)^3=28\Big(\frac{k}{2d}\Big)^3=\frac72\neq 2,\\
\E[N_4]&=4\cdot 3\Big(\frac{k}{2d}\Big)^4 +8\cdot 3\cdot 2 \Big(\frac{k}{2d}\Big)^4=\frac{15}{4}\neq 2. {\color{red} \text{this may be incorrect}}
\ea \ee  
As $n$ grows, $\E[N_n]$ will become closer to $2$ (in the case $k=2$, or to $k(k/2)^n$ in general, see also my comment further on) I think, as the contribution of the paths from the $4$ vertices with $3>d$ many neighbors decreases relatively compared to the contribution of the other paths. \\
In general (with $d=2$ and $k\leq d$), at $\ell_1$-distance $n$ there are $4(n-1)$ vertices who have, in expectation, $k/2$ many neighbors that increase their $\ell_1$-distance w.r.t.\ their parent, and $4$ vertices with, in expectation, $3k/4$ many neighbors that increase their $\ell_1$-distance w.r.t.\ their parent.\\
As a result, $W_n$ is a sub-martingale rather than a martingale. 
}
\end{comment}
we can use independence to see that 
\begin{align*}
\E[N_{n+1}|\mathcal F_n]=kN_n/2.
\end{align*}
Since $\E[W_1]=1$, we can apply the martingale convergence theorem, to ensure the existence of a random variable $W$ with 
\begin{align*}
W_n\to W\qquad \text{almost surely}.
\end{align*}
Using the second-moment method it now suffices to show that $\limsup_{n\uparrow\infty}\E[W_n^2]<\infty$ (since by the Paley–Zygmund inequality we then have that $\P(W>0)>0$, which implies $\theta^{\rm D}(k,d)>0$).

For this, note that 
\begin{align*}
\E[N_n^2]=\sum_{s,t\in \mathcal R_n}\P(s,t\text{ open})=\sum_{s,t\in \mathcal R_n}p^{K(s,t)}q^{L(s,t)}p^{2(n-K(s,t)-L(s,t))},
\end{align*}
where $p=k/(2d)$ denotes the probability that a given edge is open, $q=k(k-1)/(2d(2d-1))$ denotes the probability that two different edges that emerge from the same vertex are open, $K(s,t)$ is the number joint edges in two paths $s$ and $t$, and $L(s,t)$ is the number of vertices $x$ in $s$ and $t$ such that $s$ and $t$ do not have a joint edges after $x$.  Now note that $q<p^2$, which implies that
\[ \E[N_n^2] \leq \sum_{s,t\in \mathcal R_n} p^{2n-K(s,t)}. \]
The right-hand side is precisely the same as the right-hand side of the display below \cite[Equation~(2$\cdot$7)]{CoxDurrett}. Thus, a verbatim application of the arguments of \cite[Section 2]{CoxDurrett} implies that 
\[ \lim_{n\to\infty} \E W_n^2 <\infty \]
holds if and only if 
\[ p =k/(2d)> \varrho(d), \]
where for two independent (simple, symmetric, nearest neighbor) random walks $\widetilde S=(\widetilde S_n)_{n\in\N_0}$ and $\widetilde S'=(\widetilde S'_n)_{n\in\N_0}$ on the first quadrant of $\Z^d$ started from the origin, we define
\[ \varrho(d) = \P(\exists m\geq 0 \colon \widetilde S_m = \widetilde S'_m \text{ and } \widetilde S_{m+1} = \widetilde S'_{m+1}). \]
Let further $\tau_d=\inf \{ m \geq 0 \colon \widetilde S'_m \text{ and } \widetilde S_{m+1} = \widetilde S'_{m+1} \}$.
According to \cite[p.~155]{CoxDurrett} we have for all $d$ that
\begin{equation}\label{firstfew} 
    \P(\tau_d=0)=d^{-1}, \qquad \P(\tau_d=1)=0, \qquad \P(\tau_d=2)=d^{-3}-d^{-4}, 
\end{equation}
and for $3 \leq k \leq d$,
\[ \P(\tau=k) \leq d^{-k} k!, \]
while for $j\geq 1$ and $k >jd$,
\begin{equation}\label{bigk} 
    \P(\tau=k) \leq d^{-1} (2\pi d)^{1/2} \Big( \frac{\e^{-1/13}}{\sqrt{2\pi}} \Big)^d j^{-\frac{1-d}{2}},
\end{equation}
so that
\begin{equation}\label{CDUB} 
    \begin{aligned} 
        \varrho(d)=\P(\tau_d<\infty) &\leq d^{-1} + d^{-3} -d^{-4} + \sum_{k=3}^d d^{-k} k! +  (2\pi d)^{1/2} \Big( \frac{\e^{-1/13}}{\sqrt{2\pi}} \Big)^d \sum_{j=1}^{\infty} j^{-\frac{1-d}{2}} \\ & = d^{-1} + d^{-3} -d^{-4} + \sum_{k=3}^d d^{-k} k! +  (2\pi d)^{1/2} \Big( \frac{\e^{-1/13}}{\sqrt{2\pi}} \Big)^d \zeta\Big( \frac{d-1}{2} \Big), 
\end{aligned}\end{equation}
where $\zeta$ denotes the Riemann zeta function.
% (Note that since $\varrho(d)>1/d$, this approach cannot give any positive percolation result for $k=2$.) \color{red} András: Can we remove this sentence? Maybe it was written by myself, but it comes again in a remark below the proof, which I find more appropriate. \color{black}

%Therefore, it suffices to find the smallest $d$ such that the upper bound of~\eqref{CDUB} is less than $3/(2d)$, and this will imply that $\theta^{\rm D}(k',d') >0$ for all $k' \geq k$ and $d' \geq 3$.
%Note that this will not give any positive percolation result for $k=2$ because $k/2d = 1/d$ is smaller than the right-hand side of~\eqref{CDUB}. The next table shows the values of the right-hand side for $d=4,5,6,7$ and the respective smallest values of $k$ such that these bounds imply $\theta^{\rm D}(k,d)>0$.

The Table~\ref{table-CD} shows the smallest values of $k$ such that the right-hand side of~\eqref{CDUB} (which we computed numerically) is less than $k/(2d)$, for $d=4,5,6,7$.
\begin{table}
\begin{center}
    \begin{tabular}{|c|c|c|c|c|} \hline
       &  $d=4$ & $d=5$ & $d=6$ & $d=7$ \\ \hline
        Value of the r.h.s.\ of~\eqref{CDUB}   & 0.693093  & 0.394622  & 0.268615 & 0.199707  \\ \hline
         Smallest $k$ such that r.h.s.\ of~\eqref{CDUB} is less than $\tfrac{k}{2d} $ & 6 & 4 & 4 & 3 \\ \hline 
    \end{tabular}
\end{center}
\caption{Values of the right-hand side of~\eqref{CDUB} and the smallest $k$ such that this right-hand side is below $k/(2d) $, for $4 \leq d \leq 7$.}\label{table-CD}
\end{table}
%Via additional numerical computations we see that the right-hand side of~\eqref{CDUB} minus $k/(2d)$ is negative for any $k\geq 3$ if $7 \leq d \leq 10$ and for any $k\geq 4$ if $5 \leq d \leq 10$. 
Now, it is elementary to check that the right-hand side of~\eqref{CDUB} minus $k/(2d)$ stays negative for any $k \geq 3$ if $d \geq 7$, thanks to the following lemma:
\begin{lemma}\label{lemma-largedmon}
For any $d \geq 7$ and $k \geq 3$, we have
\begin{equation}\label{largedmon} \frac{1}{d} + \frac{1}{d^3} -\frac{1}{d^4} + \sum_{k=3}^d d^{-k} k! +  (2\pi d)^{1/2} \Big( \frac{\e^{-1/13}}{\sqrt{2\pi}} \Big)^d \zeta\Big( \frac{d-1}{2} \Big) - \frac{k}{2d} < 0.  \end{equation}
%\basc{We need a strict inequality, in the above. Do we get this from this lemma, too?}
\end{lemma}
The proof of this lemma is carried out in the Appendix~\ref{sec:largedmon}. Given that, according to Table~\ref{table-CD}, \eqref{largedmon} also holds for $k=4$ and $d=5,6$, as well as for $k=6$ and $d=4$, Lemma~\ref{lemma-largedmon} has the following consequences.  % \color{red} OK? For this it suffices to show that the rhs of~\eqref{CDUB} for dimension $d+1$ is at most $d/(d+1)$ times the same for dimension $d$, which I think is possible if $d \geq 10$. We just have to compensate the additional term $(d+1)^{-(d+1)} (d+1)!$ in the sum in dimension $d+1$. \color{black} 
First, let $k=5$. Then,  $\theta^{\rm D}(5,d) >0$ for all $d \geq 5$. Moreover, Lemma~\ref{lem:sureperc} implies that $\theta^{\rm D}(5,4)=1$. Thus, we have $\theta^{\rm D}(k,d) >0$ for $k\geq 5, d\geq 4$ (given that it is well-defined, i.e., $k\leq 2d$). For $k=4$, the table also tells us that  $\theta^{\rm D}(4,d)>0$ for all $d \geq 5$. Finally, for $k=3$, $\theta^{\rm D}(3,d)>0$ for $d \geq 7$ can be seen in the table. Therefore, to finish the proof of the theorem, we have to treat three special cases separately: $(k,d)=(4,4)$, $(k,d)=(3,5)$ and $(k,d)=(3,6)$. 

As for the first one, we compute (using a computer) the precise values of $\P(\tau_4=k)$ for $k=3,4,5$ via generating all length-6 random walk trajectories and computing how many of them belong to these events respectively (for $\P(\tau_4=\ell)$, $\ell=0,1,2$, we of course obtain the same precise values as in~\eqref{firstfew}). This yields
\[ \P(\tau_4=3)=\frac{3}{512}, \qquad \P(\tau_4=4)=\frac{279}{65\ 536}, \qquad \P(\tau_4=5)=\frac{831}{262\ 144}. \]
Thus, using~\eqref{firstfew} and~\eqref{bigk},
\[ { \begin{aligned} \P(\tau_4<\infty) & \leq \frac{1}{4} + \frac{3}{256} + \frac{3}{512} + \frac{279}{65\ 536} + \frac{831}{262\ 144} + \frac{3}{4} \sqrt{8\pi} \Big( \frac{\e^{-1/13}}{\sqrt{2\pi}} \Big)^4 + \sqrt{8\pi} \Big( \frac{\e^{-1/13}}{\sqrt{2\pi}} \Big)^4 \sum_{j=2}^{\infty} \frac{1}{j^{3/2}} \\ & \approx 0.495542 < 0.5 = \frac{k}{2d}. \end{aligned} }\]
It follows that $\theta^{\rm D}(4,4)>0$.%\basc{The first equation is incorrect as it contains an extra summation symbol, please check.} %\color{red} András: Thanks for noticing it, actually the sum from $j=1$ to $\infty$ was wrong and the one from 2 to $\infty$ was right (the term for $j=2$ we only have $3/4$ times because the $k=5$ term has been computed exactly). \color{black}

With similar methods we can show that $\theta^{\rm D}(3,5)>0$ and $\theta^{\rm D}(3,6)>0$. Indeed, for $d=5$ we obtain via explicit computer computation
\[ \P(\tau_5=3)=\frac{44}{15\ 625}, \qquad \P(\tau_5=4)=\frac{712}{390\  625}, \qquad \P(\tau_5=5)=\frac{12\ 136}{9\ 765\ 625}. \]
This already implies that  $\theta^{\rm D}(3,5)>0$ because
\[ \P(\tau_5<\infty) \leq \frac{1}{5} + \frac{4}{625} + \frac{44}{15\ 625} + \frac{712}{390\  625} + \frac{12\ 136}{9\ 765\ 625} + \sqrt{10\pi} \Big( \frac{\e^{-1/13}}{\sqrt{2\pi}} \Big)^5 \zeta(2) \approx 0.275703, \]
%\color{red} (for fun we could write that the last summand is $\sqrt 5/(24 \e^{5/13})$) \color{black} 
which is less than $k/(2d)=0.3$. For $d=6$, determining the value of $\P(\tau_6=k)$, $k=3,\ldots,6$ numerically would be computationally quite demanding. However, we see from the table that the right-hand side of~\eqref{CDUB} is just slightly larger than $k/(2d)=1/4=0.25$ when $k=3$. It turns out that in this case, it suffices to compute $\P(\tau_6=3)$ numerically, which equals $35/23328$. Estimating $\P(\tau_6>3)$ according to~\eqref{CDUB}, we obtain
\[ \P(\tau_6<\infty) \leq \frac{1}{6} + \frac{5}{1296} + \frac{35}{23\ 328} + \sum_{k=4}^6 6^{-k} k! + \sqrt{12\pi} \Big( \frac{\e^{-1/13}}{\sqrt{2\pi}} \Big)^6 \zeta(5/2) \approx 0.242338 < 0.25=\frac{k}{2d} \]
for $d=6$ and any $k\geq 3$. This resolves the three open cases, and concludes the proof. 
\end{proof}
\begin{remark}
By~\eqref{firstfew} we directly have that $\varrho(d) >1/d$, so that $\varrho(d)<k/(2d)$ never holds for $k=2$, whence for $k=2$ the approach of the proof of Theorem~\ref{thm-CoxDurrett} is not applicable. For $d\leq 3$ (where the cases $k=2,d=2$, $k=2,d=3$, and the case $k=3,d=3$ are open), the issue is that the sum $\sum_{j=1}^{\infty} j^{\frac{-1-d}{2}}$ (cf.\ \eqref{CDUB}) does not converge. Indeed, this technique, like many others in statistical physics, only works when we are in at least four dimensions.
Finally, for $d=4$ and $k=3$ one could hope that a combination of the proof techniques of the theorem and some explicit computations can work, but for this, one would need a sufficiently tight upper bound on $\P(\tau_4 = k)$ up to $k \approx 12$, which exceeds our available computing capacity. 
\end{remark}

The following theorem provides a coupling between the $k$-DnG in $d$ dimensions and the $(k+1)$-DnG in $d+1$ dimensions and, in particular, provides a relation between the percolation probabilities in these two models.

\begin{theorem}\label{theorem-3322}
For all $k,d\ge1$ we have that $\theta^{\rm D}(k+1,d+1)\ge \theta^{\rm D}(k,d)$.
\end{theorem}
\begin{proof}
We use a randomized coupling approach. The coupling is designed to have the following two key properties. First, using the additional dimension, the edge distribution of a node in $\Z^d$ is the image measure of a mapping from the nodes "above and below" that node and hence i.i.d.\ over the nodes in $\Z^d$. Second, any connected component in $\Z^d$ is the image of a connected component in $\Z^{d+1}$.

To make this precise, let us write $x=(x_1,\dots,x_{d+1})\in \Z^{d+1}$ and $x'=\pi(x)=(x_1,\dots, x_d)\in \Z^d$ for the projection of $x$ to its first $d$ coordinates. By $\omega_{x}$ we denote the $(d+1)$-dimensional $(k+1)$-neighbor directed edge configuration of $x\in \Z^{d+1}$ and by $\omega'_{x'}=\pi(\omega_{x})$ the associated edge configuration of $x'$ in $\Z^d$. In words we forget all edges facing up or down and only keep all remaining edges which we consider as part of the edge set of $\Z^d$. 

We further write $\bar x=(x',n)_{n\in\Z}$ for the vector of nodes {\em above and below} $x'\in \Z^d$ in the $(d+1)$-th dimension and $\bar\omega_{\bar x}$ for the associated vector of $(k+1)$-neighbor configurations of the nodes $(x',n)\in \Z^{d+1}$ in $\bar x$. The idea now is to define probability kernels $\nu_{x'}(\cdot|\bar\omega_{\bar x})$, step-by-step with respect to $x'$, in such a way that the joint kernel has two properties:
\begin{enumerate}
\item Integrating the joint kernel with respect to $\omega$ gives the distribution of the $d$-dimensional directed $k$-DnG. 
\item If the connected component of the origin in a $(d+1)$-dimensional component $\omega$ is finite, then the joint kernel only puts positive mass on finite components of the origin in $d$ dimensions as well. 
\end{enumerate}

To define this precisely, we need to consider ordered $k$-subsets of nearest neighbors of nodes $x'\in \Z^d$ with some additional coordinate, i.e., we write 
$$D_{x'}=\{\{(a_{1},n_1),\dots,(a_{k},n_k)\}\colon n_1,\dots,n_k\in \Z, \{a_1,\dots,a_k\}\subset N^d(x')\text{ with }a_i\neq a_j\text{ for all }i\neq j\},$$
where $N^d(x')\subset\Z^d$ denotes the set of nearest neighbors of $x'$ in $\Z^d$.
Then, we start with the origin $o\in \Z^d$ and define for $\bar\omega_{\bar o}$ the probability kernel $\nu_{o,0}(\cdot|\bar\omega_{\bar o})$ as a probability measure on $D_o$ as follows. 
\begin{enumerate}
\item[\bf Step 1] If $|\pi(\omega_{(o,0)})|\ge k$, pick $k$ neighbors $a_{o,1},\dots,a_{o,k}$ uniformly at random from the available $|\pi(\omega_{(o,0)})|$ connected neighbors in $\Z^d$. Write the outcome as $\big\{(a_{o,1},0),\dots,(a_{o,k},0)\big\}$.
\item[\bf Step 2] If $|\pi(\omega_{(o,0)})|<k$, we keep the $(k-1)$-many available connected neighbors in $\Z^d$ as \\ $a_{o,1},\dots, a_{o,k-1}$. Then, randomly choose a direction up or down (there must exist a directed edge facing up as well as down) and follow these steps:
\begin{enumerate}
\item If $|\pi(\omega_{(o,\pm 1)})\setminus \pi(\omega_{(o,0)})|\ge 1$, 
pick the missing neighbor $a_{o,k}$ uniformly at random from the available connected neighbors $\pi(\omega_{(o,\pm 1)})\setminus \pi(\omega_{(o,0)})$. Write for the outcome $\big\{(a_{o,1},0),\dots,(a_{o,k-1},0),(a_{o,k},\pm 1)\big\}$.
\item If $|\pi(\omega_{(o,\pm 1)})\setminus \pi(\omega_{(o,0)})|= 0$, follow the existing arrow in the same direction as in Step~2 and repeat from (a) with $\omega_{(o,\pm 1)}$ replaced by $\omega_{(o,\pm 2)}$.
\end{enumerate}
\end{enumerate}
Now, almost surely, the construction ends and, due to the complete symmetry in the construction, $\nu_{o,0}$ reproduces the uniform $k$-nearest-neighbor distribution for the origin. That is, for 
$$E^{k,d}_{x'}=\{\{a_1,\dots, a_k\}\subset N^d(x')\}$$ 
and $F\colon D_{x'}\to E^{k,d}_{x'}$, $\big\{(a_1,n_1),\dots,(a_k,n_k)\big\}\mapsto \{a_1,\dots, a_k\}$ and $U^{k,d}_x$ the uniform distribution on $E^{k,d}_x$, we have that, for all $\omega'_o\in E^{k,d}_o$,
\begin{align*}
\int \Big(\bigotimes_{n\in \Z}U^{k+1,d+1}_{(o,n)}\Big)(\d \bar \omega_{\bar o})\nu_{o,0}(F^{-1}(\omega'_o)|\bar\omega_{\bar o})=U^{k,d}_o(\omega'_o).
\end{align*}

Let us describe the following steps in words. Under $\nu_{o,0}(\cdot|\bar\omega_{\bar o})$ we are equipped with $k$ neighbors of the origin in $\Z^d$ that all also carry the information on which level in the additional $(d+1)$-th coordinate they are discovered. Starting with $\hat a_{o,1}=(a_{o,1},n_1)$, the first discovered neighbor, we can sample its neighbors, using the same algorithm, based on the information provided by $\bar a_{o,1}$, now started at level $n_1$ (and not at level $0$). The same can be done for all other neighbors $\hat a_{o,i}=(a_{o,i},n_i)$. In this fashion, we can step-by-step explore the connected component of the origin in $\Z^d$, where in any step, whenever we discover a new vertex $y\in \Z^d$, we only use information from the associated vector $\bar y$, i.e., 
\begin{align*}
\nu(\d \omega'|\omega)=\int &\nu_{(o,0)}(\d(\hat a_{o,1},\dots,\hat a_{o,k})|\bar \omega_{\bar o})\times\\
&\int\nu_{(a_{o,1},n_1)}(\d(\hat a_{a_{o,1},1},\dots,\hat a_{a_{o,1},k})|\bar \omega_{\bar a_{o,1}})\dots \int\nu_{(a_{o,k},n_k)}(\d(\hat a_{a_{o,k},1},\dots,\hat a_{a_{o,k},k})|\bar \omega_{\bar a_{o,k}})\times\\
&\dots \1\{F(\hat a_{o,1},\dots,\hat a_{o,k},\hat a_{a_{o,1},1},\dots,\hat a_{a_{o,1},k},\dots)=\d \omega'\}.
\end{align*}

Now, the key point is the following. Imagine that the algorithm presents us a configuration of open edges in $\Z^d$ such that the connected component of the origin reaches the boundary of a centered box with side-length $n$. Then, necessarily, any configuration in $\Z^{d+1}$ that is used as an input for the algorithm must also have the property that the origin is connected to the boundary of a box of side-length $n$, now of course in  $\Z^{d+1}$. Hence, denoting by $A^{d}_n$ the event that the origin is connected to the boundary of the centered box of side-length $n$ in $\Z^d$ and by $\P_{k,d}$ the distribution of the $k$-DnG in dimension $d$, we have that 
\begin{align*}
\P_{k,d}(A^d_n)=\E_{k+1,d+1}[\nu(A^d_n|\cdot)]\le \P_{k+1,d+1}(A^{d+1}_n), 
\end{align*}
which gives the result.
% Further let $E^i$ denote the edge set associated to $\Z^i$ for $i\in\{2,3\}$. 
% $\omega$ we will denote a configuration of open directed edges in $E^3$ and by
% $E^3_{\bar x}$ 
% $E^2_x$ denote the set of two-dimensional two-neighbor directed edge configurations $\omega_x$ of $x$. Correspondingly, 
% Let $\pi$ denote the projection of the edges from $\omega_{(x,n)}$ to the edges of $x$ in $\Z^2$, in words we forget all edges facing up or down and only keep all remaining edges which we consider as part the edge set of $\Z^2$. 
% We now define, in an algorithmic fashion, a probability kernel from $(E^3_{\bar x})$
% Let us start by  defining a probability kernel $\nu_x$ from $E^3_{\bar x}$ to probabilities on $E^2_x$. 
% \begin{enumerate}
% \item If $|\pi(\omega_{(x,0)})|=3$, pick two of the three directed edges uniformly at random and define the result as $\omega_x=F(\pi(\omega_{(x,0)}))$.
% \item If $|\pi(\omega_{(x,0)})|=2$, define $\omega_x=\pi(\omega_{(x,0)})$.
% \item If $|\pi(\omega_{(x,0)})|=1$, keep the edge in $\Z^2$ and follow the existing arrow upwards, consider $\omega_{(x,1)}$ and do the following:
% \begin{enumerate}
% \item If $|\pi(\omega_{(x,1)})|\ge 2$, pick one of the directed edges in $\pi(\omega_{(x,1)})\cap \pi(\omega_{(x,0)})$ uniformly at random and define the result as $\omega_x=F(\pi(\omega_{(x,0)}),\pi(\omega_{(x,1)}))$.
% \item If $|\pi(\omega_{(x,1)})|=1$ and $\pi(\omega_{(x,1)})\neq \pi(\omega_{(x,0)})$, define $\omega_x=\pi(\omega_{(x,0)})\cup\pi(\omega_{(x,1)})$.
% \item If $|\pi(\omega_{(x,1)})|=1$ and $\pi(\omega_{(x,1)})= \pi(\omega_{(x,0)})$, follow the existing arrow upwards and repeat from (a) with $\omega_{(x,1)}$ replaced by $\omega_{(x,2)}$.
% \end{enumerate}
% \end{enumerate}
% Now, almost surely, for every $x\in\Z^d$, the construction ends. Further, due to the complete symmetry in the construction, the resulting field $(\omega_x)_{x\in\Z^2}$ has the same distribution as the directed 2NN graph, i.e., $$\E_3[\nu(\d\omega|\cdot)]=\P_3(\d\omega).$$
% The essential benefit of the coupling comes from the following fact. Let  $\bar\omega$ be such that the origin is not connected to the boundary of a box with sidelength $n$, i.e., $\bar\omega\notin A^3_n$, then also $\nu(A^2_n|\bar\omega)=0$, since we are only using edges that are already present in $\bar\omega$. THIS I GUESS IS NOT QUITE CORRECT YET; WE HAVE TO DESIGN THE ALGORITHM STARTING FROM THE ORIGIN. 
% But then, 
% \begin{align*}
% \P_3(A^2_n)=\E_3[\nu(A^2_n|\cdot)]\le \P_3(A^3_n), 
% \end{align*}
% which tends to zero as $n$ tends to infinity, by assumption.
\end{proof}

We see from the proof of Theorem~\ref{theorem-3322} that it is essential that \emph{both} $k$ and $d$ increase by one to make the coupling work. For example, issues arise when trying to deduce a similar coupling between the settings $k=d=2$ and $k=2,d=3$. Though we expect percolation occurs in both settings, a similar coupling would give rise to a two-dimensional model where some vertices have only out-degree one. In such a case, it seems that the additional edges that appear in the $2$-DnG in two dimensions are pivotal in the sense that they enable percolation to occur whereas it may not occur without the presence of these edges. 

Still, we expect that a relation between the $d$-dimensional and $(d+1)$-dimensional settings should exist, and thus state the following conjecture.

\begin{conj}\label{conj:dimensioncouple}
For all $k,d\ge 2$ such that $k\leq d$, we have that $\theta^{\rm D}(k,d+1)\geq \theta^{\rm D}(k,d)$. 
\end{conj}

The intuition behind this statement is as follows: The out-degree of each vertex is the same in both models. Hence, if we were to explore the connected component of the origin in a breadth-first search (BFS) manner, it should be less likely for vertices in the $n$-th generation of the BFS exploration to discover common new neighbors when the dimension increases, for any $n\in\N$. Thus, the sizes of the generations of the BFS exploration in $d+1$ dimensions should dominate those in $d$ dimensions. In essence, the grid $\Z^d$ becomes more `tree-like' and the BFS exploration is `better approximated' by a branching process with mean $k$ offspring when $d$ increases.

The condition $k\leq d$ follows from Lemma~\ref{lem:sureperc}, since $\theta^{\rm D}(k,d)=1$ whenever $k\geq d+1$. Hence, the cases $k\geq d+2$ are trivially true by this lemma, since then $\theta^{\rm D}(d+2,d)=\theta^{\rm D}(d+2,d+1)=1$, whilst for $k=d+1$ it may be true that $\theta^{\rm D}(d+1,d+1)<\theta^{\rm D}(d+1,d)=1$.

% It would be desirable to proof the following conjecture.
% \begin{conj}
% $\theta^{\rm D}(2,2)>0$. 
% \end{conj}
% \basc{Do we want to move this to an open problem section?}


\subsection{Results for the undirected $k$-neighbor graph}\label{sec:U}
% \begin{conj}
% We have that $\theta^{\rm U}(k,d)>0$ for all $k,d\ge 2$.  
% \end{conj}
% \basc{Again, move this to open problems?} \color{red} András: Yes, I agree in both cases. I would just remove the assertions and refer the reader (perhaps already before the beginning of Section 2.1) to the open problem section. But if you guys like it better with open problems separated according to the kind of graph, we can stick to it. \color{black}

\medskip

It is clear that $\theta^{\rm U}(k,d) \geq \theta^{\rm D}(k,d)$ for any $k,d$ because any directed edge in the $k$-DnG corresponds to an undirected edge in the $k$-UnG (and at most two directed edges can correspond to the same undirected edge). Additionally, $k=1$ is not sufficient for percolation even in the undirected sense.  

\begin{prop}\label{prop-U1}
    We have that $\theta^{\rm U}(1,d)=0$ for all $d \geq 1$.
\end{prop}

For the proof of Proposition~\ref{prop-U1} as well as some further results, we introduce the following notation. Let us denote by $c(d)$ the connective constant of $\Z^d$ (cf.~e.g.\ \cite{Grimmett1999}), that is $$c(d)=\lim_{n\uparrow\infty}c_n(d)^{1/n},$$
where $c_n(d)$ is the number of self-avoiding paths in the $d$-dimensional cubic lattice, starting at the origin, of length $n$. The limit exists by subadditivity. Also, $d<c(d)<2d-1$. 

\begin{proof}
Assume for a contradiction that the assertion of the proposition is wrong, and let $A$ then denote the event of positive probability that there exists an infinite path starting from $o$ consisting of edges of the $1$-UnG. %First, we claim that if such a path exists, it can contain at most one bidirectional edge. Indeed, two bidirectional edges cannot be adjacent to each other because this would mean that the vertex separating them has outdegree at least 2 in the $1$-DnG. Also, vertices belonging to the bidirectional edge have no outgoing edge but the one contributing to the bidirectional edge. This way, if $Y_0=o, Y_1,Y_2,\ldots$ is the enumeration of the vertices among an infinite path and for some $i \geq 0$, $(Y_i,Y_{i+1})$ is a bidirectional edge, then either $Y_{i+2}=Y_i$ and $Y_i$ has an additional incoming edge in the $1$-DnG coming from $Y_{i+3} \neq Y_{i+1}$, or $Y_{i+2} \neq Y_i$ and $Y_{i+1}$ has an additional incoming edge coming from $Y_{i+2}$. Now, 
%such a path, then if there is a bidirectional edge, say $(Y_{i},Y_{i+1})$, then the only way that this edge is in the infinite path is that $Y_i$ or $Y_{i+1}$ has an additional incoming edge from some other vertex
%there also exists one with the same property that does not contain any bidirectional edge apart from the very first one (which is then adjacent to $o$). Indeed, in any infinite path of the $1$-UnG, there cannot be two bidirectional edges adjacent to each other because this would mean that the vertex between them has outdegree at least 2. Thus, taking any infinite path starting from $o$, this can only contain isolated bidirectional edges (i.e., ones not adjacent to any other one). Hence, if $Y_0=o,Y_1,Y_2,\ldots$ is an infinite path of the $1$-UnG and for some $k \geq 0$, $(Y_k,Y_{k+1})$ is a bidirectional edge, it must be the case that either $k=0$, or $(Y_{k-1},Y_k)$ and $(Y_{k+2},Y_{k+1})$ is an edge of the $1$-DnG. In either of the cases, if we leave out the bidirectional edge $(Y_k,Y_{k+1})$ from the infinite path, the path still remains infinite, and this can be done simultaneously with all bidirectional edges of the path. This implies the claim.
Let $o=X_0,X_1,X_2,\ldots$ be any infinite path of the $1$-UnG starting from $o$.
%that contains no bidirectional edges apart from $(X_0,X_1)$ and is defined (on the event $A$) as follows. First, consider the infinite cluster of $o$ in the $1$-UnG and remove all bidirectional edges of it that are not adjacent to $o$, along with the vertices that are only adjacent to a bidirectional edge of the cluster. Now, we define $X_1,X_2,\ldots$ recursively as follows. If $o$ is not adjacent to a bidirectional edge of the infinite cluster, then given that $X_0,\ldots,X_{i-1}$ is already defined, let $X_i$ be the lexicographically smallest one among all vertices $v$ of the remainder of the infinite cluster of $o$ such that $(X_{i-1},v)$ is an edge of the $1$-UnG and there is a path in the $1$-UnG from $v$ to infinity that does not use any of the edges $(X_0,X_1),\ldots,(X_{i-1},v)$. (Note that there may be several $v$ satisfying this property in case there are multiple edge-disjoint incoming paths of the $1$-DnG at $X_{i-1}$.) If $o$ is adjacent to a bidirectional edge, then let $X_1$ be its neighbor via this edge, and define $X_2,X_3,\ldots$ as before.

Now, if $X_0,X_1,X_2$ is the sequence of vertices in an infinite path (on the event $A$), we put
\[ K=\inf \{ k \geq 0 \colon (X_k,X_{k+1}) \text{ is not an edge of the $1$-DnG} \}. \]
We know from (the proof of) Lemma~\ref{lem:surenoperc} that $K<\infty$ almost surely. Next, let us define
\[ L=\inf \{ k > K \colon 
(X_{k},X_{k+1})\text{ is an edge of the 1-DnG}\}. \]
Now we claim that $L=\infty$ on the event $\{ K<\infty\}$ (where we put $K=\infty$ on the event $A^c$). Indeed, for all $k=K,K+1,\ldots,L-1$, $(X_{k+1},X_k)$ is an edge of the $1$-DnG because it follows from the definition of $K$ that $(X_{k},X_{k+1})$ is not an edge of the 1-DnG, but $(X_k,X_{k+1})$ is an edge of the 1-UnG, and this is only possible if $(X_{k+1},X_k)$ is an edge of the $1$-DnG. But now, if $L<\infty$, then $(X_{L},X_{L-1})$ is an edge of the $1$-DnG, so that $(X_{L},X_{L+1})$ cannot be an edge of the $1$-DnG since there is only one edge going out of $X_L$. This implies the claim.

%\color{blue} Note that this implies that the cluster of $o$ in the $1$-UnG consists of coalescing (finite or infinite) directed paths towards $X_K$ and $X_{K+1}$ if $K\geq 1$; if $K=0$, then all are directed towards $o$. In particular, there can be at most one bidirectional edge in the infinite cluster, namely the one between $X_{K-1}$ and $X_K$ (because if two vertices have outgoing edges of the $1$-DnG towards each other, they cannot have any other outgoing edges, so that $K$ is at most the larger one among the indices of the two endpoints of the bidirectional edge). We perhaps don't need this, but I wanted to note it. \color{black}

Let us denote the $\ell_1$-sphere of radius $n\in\N_0$ by $S^{(1)}_n$ (with $S^{(1)}_0=\{o\}$). From Lemma~\ref{lem:surenoperc} we even know that $\P(\kappa<\infty\,|\,A)=1$, where
\[ \kappa = \inf \big\{ k \geq 0 \colon o \not\rightsquigarrow S^{(1)}_k \text{ in the $1$-DnG} \big\}. \]
Since $o \in S_0^{(1)}$ and $o$ is always connected to one of its nearest neighbors by a directed edge, we have $\P(\kappa \geq 2)=1$.
Given that $\kappa$ is an $\N_0 \cup \{\infty\}$-valued random variable, we can find $k_0\in\N \setminus \{ 1 \}$ and $\varepsilon>0$ such that $\P(\kappa =  k_0 |A) \geq \varepsilon$.  Now, by our previous observation, since from $o$ to $S^{(1)}_{\kappa-1}$ there always exists a directed path in the $1$-DnG, 
on the event $\{ \kappa = k_0 \} \cap A$, there must exist a self-avoiding directed path of length $n$ in the $1$-DnG ending at some vertex of $S^{(1)}_{k_0-1}$  for all $n\in\N_0$. %\basc{This argument is not clear to me. Isn't it the case that three \emph{doesn't} exist a path from $o$ to $S_\kappa ^{(1)}$ in the $1$-DnG? Or do we want to use $S_K^{(1)}$ here instead? No, that makes no sense either...} 
Thus, writing $c_{n,k_0-1}(d)$ for the number of self-avoiding paths in the $d$-dimensional lattice of length $n$ starting from (or ending at) $S^{(1)}_{k_0-1}$, it is clear that
\[ \lim_{n\to\infty} c_{n,k_0-1}(d)^{1/n} = \lim_{n\to\infty} c_{n,1}(d)^{1/n}=c(d) \leq 2d-1, \]
where $c(d)$ is the connective constant of $\Z^d$. Now, any self-avoiding path of length $n$ ending at $S^{(1)}_{k_0-1}$ is open with probability $(2d)^{-n}$. Hence,
\[ \P(\exists\text{ a self-avoiding path of length $n$ starting from $S^{(1)}_{k_0-1}$ included in the $1$-DnG}) \leq \Big( \frac{2d-1}{2d} \Big)^n, \]
which decays exponentially fast in $n$. This contradicts the assertion that the left-hand side is bounded from below by $\P(A)\varepsilon>0$ independently of $n$. Therefore, $\P(A)=0$, as desired.
\end{proof}

Next, we verify percolation in the $2$-UnG and $3$-UnG models in two and three dimensions, respectively.

\begin{theorem}\label{theorem-U22}
  We have $\theta^{\rm U}(2,2)>0$ and $\theta^{\rm U}(3,3)>0$.  
\end{theorem}

\begin{proof}
    We prove the positive percolation probability in the 2-UnG via a dual approach. That is, let $\D^2:=\{x+(1/2,1/2): x\in \Z^2\}$ be the two-dimensional dual lattice of $\Z^2$. An edge $e$ in the dual lattice has exactly one edge $e'$ from the lattice $\Z^2$ that crosses $e$, and we say the edge $e$ is closed if and only if the edge $e'$ is open. For the $2$-UnG model in two dimensions, it thus follows that for any such edge $e$, 
    \be 
    \P(e \text{ is closed})=1-\P(e'\text{ is closed})=1-3/4=1/4.
    \ee 
    Moreover, the negative correlations between edges that are present in the $2$-UnG model also appear in its dual. Namely, for two edges in the dual that are vertex disjoint or not orthogonal, the status of the edges is independent. When two edges $e_1,e_2$ in the dual share a vertex and are orthogonal, that is, when there is some vertex $v\in \Z^2$ such that the edges emanating from $v$ can cross both $e_1$ and $e_2$, then
    \be 
    \P(e_1\text{ and }e_2\text{ closed in the dual})=1/24\leq (1/4)^2=\P(e_1\text{ closed in the dual})^2.
    \ee 
    This can be extended, so that for any $k\geq 2$ and any path $(v_1,\ldots, v_k)$ of nearest-neighbor vertices in $\D^2$, 
    \be 
    \P(\text{Path }(v_1,\ldots, v_k)\text{ closed})\leq \prod_{i=1}^{k-1}\P((v_i,v_{i+1})\text{ closed})=(1/4)^{k-1}.
    \ee 
    As such, we can use the negative correlation to bound the probability that there exists a closed circuit in the dual that surrounds the origin. For any $n\geq 4$, the number of such circuits of length $n$ is at most $nc_{n-1}(2)$, where we recall that $c_n(2)$ denotes the number of self-avoiding paths in $\Z^2$. As $c_n(2)=(c(2)+o(1))^{n-1}$, we thus arrive at 
    \be \ba\label{eq:finnorpercprob}
    1-\theta^{\rm U}(2,2)&=\P(\exists \text{ circuit }\gamma \text{ in }\D^2\text{ that surrounds the origin})\\
    &\leq \sum_{n=4}^\infty \sum_{\substack{\text{circuits }\gamma\\ |\gamma|=n}}\P(\gamma\text{ closed in }\D^2)\\
    &\leq \sum_{n=4}^\infty n(c(2)+o(1))^{n-1} (1/4)^n, 
    \ea \ee 
    which is finite since $c(2)\leq 3$. We now apply an argument from~\cite{JanSwart} (based on an argument from Durrett~\cite{Durr88})  to show that this probability is in fact smaller than one, which yields $\theta^{\rm U}(2,2)>0$. Let $D_m=\{0,\ldots, m\}^2$ and say that a vertex $i\in \Z^2$ is `wet' when there exists a $j\in D_m$ such that $i\to j$, that is, when there exists a path of closed edges from $i$ to $j$. Suppose the component of the origin is finite almost surely. Then, the number of wet vertices is finite almost surely as well, and there exists a circuit in the dual that surrounds all the wet sites. By the same argument as in~\eqref{eq:finnorpercprob}, 
    \be 
    \P(D_m \not\to \infty )\leq \sum_{n=4m}^\infty \sum_{\substack{\text{circuits }\gamma \\ |\gamma|=n}}\P(\gamma\text{ closed in }\D^2)\leq \sum_{n=4m}^\infty n(c(d)+o(1))\Big(\frac14\Big)^n,
    \ee 
    where we can start the outer sum from $4m$ since a circuit that surrounds all wet vertices must surround $D_m$ and hence have length $4m$ at least. 
    As a result, since the sum is finite for each $m\in\N$, choosing $m$ large enough yields that $\P(D_m\to \infty)>0$ and hence $\P(i\to \infty)>0$ for some $i\in D_m$. By translation invariance, it then follows that $\theta^{\rm U}(2,2)>0$, as desired.

    To prove that $\theta^{\rm U}(3,3)>0$, it suffices to show that with positive probability $o$ is connected to infinity via any two-dimensional plane $S$ of $\Z^3$ including $o$. This can be done via the dual approach employed in the first part of the proof. In the restriction of the $3$-UnG to $S$, each edge of the $3$-UnG is open with probability $3/4$. Thus, for any edge $e$ of the dual lattice corresponding to $S$, we have $\P(e\text{ is closed})=1/4$. For any two dual edges that are vertex disjoint or not orthogonal, their status is independent, and if two dual edges $e_1,e_2$ share a vertex and are orthogonal, then we have
\[ \P(\text{$e_1$ and $e_2$ closed in the dual}) = \frac{\binom{5}{3}}{\binom{6}{3}} \frac{\binom{5}{3}}{\binom{6}{3}}\frac{\binom{4}{3}}{\binom{6}{3}} = \frac{1}{20}\leq \Big(\frac{1}{4}\Big)^2 = \P(e_1\text{ closed in the dual})^2, \]

whence the proof can be finished analogously to the one of Theorem~\ref{theorem-U22}.
\end{proof}

\begin{remark}
One can verify percolation in the $d$-UnG in $d$ dimensions for any $d \geq 4$ similarly to the case $d=3$ (or, in fact, for any $k\geq 2, d\geq 3$), but this is not necessary since the same assertion already follows from Theorem~\ref{thm-CoxDurrett}. 
% \basc{We can do this more generally for any combination of $k$ and $d$, not just $k=d$, similar to the proof of Proposition~\ref{prop-Bperc} This would yield the condition
% \be 
% \Big(1-\frac{k}{2d}\Big)^2<\frac{1}{c(2)}, \qquad \text{or, equivalently,}\qquad k>2d\Big(1-\frac{1}{\sqrt{c(2)}}\Big).
% \ee 
% Do we want to include this here (or perhaps as a remark after Proposition~\ref{prop-Bperc}), or is this not very interesting? It doesn't provide any new results (at least, not the open cases in Table~\ref{table-UD}) and I think the $k$ that satisfy this condition already are captured by the $k$-DnG in Theorem~\ref{thm-CoxDurrett}.
% } \color{red} András: If you also agree, Bene and Jonas, we can remove this. \color{black}
\end{remark}

\subsection{Results for the bidirectional $k$-neighbor graph}\label{sec:B}
Let us recall the connective constant $c(d)$ from Section~\ref{sec:U}. We have the following lemma.
\begin{lemma}\label{lemma-Bnoperc}
For all $k$ such that $k(k-1)<2d(2d-1)/c(d)$ we have that $\theta^{\rm B}(k,d)=0$.
\end{lemma}
Let us note that this result can be interpreted in two ways. First, for given $d$, absence of percolation is guaranteed for sufficiently small $k$, and for example $\theta^{\rm B}(2,d)=0$ for any $d\geq 2$ (while it is clear that $\theta^{\rm B}(2,1)=1$), $\theta^{\rm B}(3,d)=0$ for any $d \geq 3$ and $\theta^{\rm B}(4,d)=0$ for any $d \geq 6$, as can be seen from the lower bounds on $c(d)$ by \cite{HSS}. However, also for fixed $k$, since $2d(2d-1)/c(d)>2d$, for sufficiently large $d$ there is no percolation. This is due to the fact that in high dimensions, it is unlikely for two neighboring vertices to pick the same connecting edge. 

\begin{proof}
Note first that the system is negatively correlated in the sense that for all directed edges $\ell_1=(x_1,x_2)\neq\ell_2=(y_1,y_2)$ we have that
$$
\mathbb P(\ell_1\text{ open and }\ell_2\text{ open})\le \mathbb P(\ell_1\text{ open})^2.
$$
Indeed, note that if $x_1\neq y_1$, then the inequality is an equality by independence. However, if $x_1=y_1$ then 
$$
\mathbb P(\ell_1\text{ open}\,|\,\ell_2\text{ open})=(k-1)/(2d-1)\le k/(2d)=\mathbb P(\ell_1\text{ open}).
$$
Note that the negative correlation carries over also to the bidirectional (and also undirected) case since they are build from the directed case. Let us denote by $\ell_n$ the straight line starting at the origin up to the node $(ne_1)$. Then, using the first moment method and negative correlations, we have that 
\begin{align*}
\P(o\leftrightsquigarrow_{\rm B}\partial B_n)\le \sum_{s\in \mathcal R_n}\P(s\text{ is open})\le c(d)^n\P(\ell_n\text{ is open}).
\end{align*}
Finally, for a suitable constant $C>0$, we have that $\P(\ell_n\text{ is open})\le C p^n$ where $p=k(k-1)/(2d(2d-1))$ is the probability that the origin chooses precisely two prescribed edges. This leads to the criterion for absence of percolation $c(d)p<1$.
\end{proof}

The approach of verifying percolation restricted to a two-dimensional plane, which we used in order to derive $\theta^{\rm U}(3,3)>0$ in the proof of Theorem~\ref{theorem-U22}, is also applicable in the bidirectional case, as the following proposition shows.
\begin{prop}\label{prop-Bperc}
We have $\theta^{\rm B}(k,d) > 0$ whenever 
\be 
 k>d\sqrt{4\Big(1-\frac{1}{c(2)}\Big)}.
\ee 
\end{prop}

An application of the upper bound $c(2)\leq 2.679192495$ from~\cite{PonTitt00} for the connective constant of $\Z^2$ immediately yields the following corollary. 

\begin{corollary}\label{cor-Bperc}
We have $\theta^{\rm B}(k,d) > 0$ whenever 
\be 
 k>d\sqrt{4\Big(1-\frac{1}{2.679192495}\Big)}\approx 1.583355
d.
\ee 
\end{corollary}
%\basc{Do we want to use this corollary, or just put the best known bound for $c(2)$ in the proposition straight away?} \color{red} András: I don't mind having it as a separate corollary. An argument against it could be that the bound 2.6791... is very close to the true value of $c(2)$ and there's not much space for improvement. \color{black}
%\basc{This small improvement, just using a better bound than $3$ on the connective constant by~\cite{PonTitt00} improves percolation for $(2d-2)$-BnG from $d\geq 6$ to $d\geq 5$, $(2d-3)$-BnG from $d\geq 9$ to $d\geq 8$, and $(2d-4)$-BnG from $d\geq 11$ to $d\geq 10$.}

\begin{remark}
This corollary allows us to verify percolation e.g.\ for the $(2d-1)$-BnG for $d \geq 3$, for the $(2d-2)$-BnG for $d \geq 5$, for the $(2d-3)$-BnG for $d \geq 8$, and for the $(2d-4)$-BnG for $d \geq 10$. (Thus, the smallest-dimensional positive percolation result that we obtain is that $\theta^{\rm B}(5,3) >0$.)
\end{remark}

\begin{proof}[Proof of Proposition~\ref{prop-Bperc}]
As in the proof of the assertion $\theta^{\rm U}(3,3)>0$ of Theorem~\ref{theorem-U22}, it suffices to verify that $o \rightsquigarrow \infty$ in the restriction of the $k$-BnG to a fixed two-dimensional plane $S$ including $o$ with positive probability, using the duality approach of the proof of Theorem~\ref{theorem-U22}. Dual edges of the two-dimensional lattice of $S$ are open with probability  $1-(k/2d)^2$. %\color{black} \basc{Please check, I think this is not correct. First, the edge being open means that it is not present, right? The probability that the dual edge is not present, equals the probability that the edge in the $k$-BnG is present, which equals 
%\be 
%\P(\text{Edge present in $k$-BnG})=\P(\text{Both directed edges present in $k$-BnG})=\bigg(\frac{\binom{2d-1}{k-1}}{\binom{2d}{k}}\bigg)^2=\Big(\frac{k}{2d}\Big)^2.
%\ee 
%The remainder of the argument in the proof is correct, however.}
Further, the indicators of two dual edges being closed, i.e., the corresponding edges of the $k$-BnG restricted to $S$ being open, are non-positively correlated. Indeed, we saw already in the proof of Proposition~\ref{lemma-Bnoperc} that the indicators of two edges of the $k$-BnG being open are either independent or strictly negatively correlated. Now, if the indicators of two events are independent or strictly negatively correlated, then so are the indicators of the complements of the two events. We see from the proof of Theorem~\ref{theorem-U22} that given these non-positive correlations, it suffices to choose $k$ and $d$ in such a way that a dual edge is closed with probability at most $1/c(2)$. This holds whenever 
\be  
\Big(\frac{k}{2d}\Big)^2 > 1-\frac{1}{c(2)},\qquad \text{or, equivalently,}\qquad k>d\sqrt{4\Big(1-\frac{1}{c(2)}\Big)}, 
\ee
as asserted. 
\end{proof}

Let us now present a statement on the existence of a supercritical phase in high dimensions.

\begin{prop}\label{prop-1indep}
%\begin{enumerate}[(i)]
%\item For any $d \geq 2$, if $k>2\sqrt{0.8457}d \approx 1.8392d$, then $\theta^{\rm B}(k,d)>0$.
%\item 
For any $\alpha > 2 \sqrt{0.5847} \approx 1.5293$, we have $\theta^{\rm B}(\alpha d, d)>0$ for all $d$ sufficiently large.
%\end{enumerate}
\end{prop}

\begin{proof}
Note that the $k$-BnG model (just as the $k$-UnG model) in any dimension $d$ is an 1-dependent Bernoulli bond percolation model where each edge is open with probability $p=k^2/(4d^2)$. Let us write 
\[ p_{\max}(\Z^d)= \sup \{ p \in (0,1) \colon \text{some 1-dependent Bernoulli bond perc.\ model does not percolate} \}. \]
According to \cite{improvedLB}, we have %$p_{\max}(\Z^d) \leq p_{\max}(\Z^{d-1})$, %and $p_{\max}(\Z^2) \leq 0.8457$. These assertions imply (i). Part (ii) follows from the assertion of the same paper that 
further, $\lim_{d\to\infty} p_{\max}(\Z^d) \leq 0.5847$, which implies the proposition.
\end{proof}
\begin{remark}
Another result by \cite{improvedLB} is that $p_{\max}(\Z^2) \leq 0.8457$. This together with the assertion that $p_{\max}(\Z^d) \leq p_{\max}(\Z^{d-1})$ yields that for $\alpha>2\sqrt{0.8457}$ we have $\theta^{\rm B}(\alpha d, d)>0$. However, this assertion is weaker than the one of Proposition~\ref{prop-Bperc} because $2\sqrt{0.8457} \approx 1.8392 > \sqrt{4(1-1/2.679192495)}$. 
%The smallest nontrivial (i.e., such that $k \neq 2d$) case covered by Proposition~\ref{prop-1indep} is $k=13$ and $d=7$: Since $13/7 >1.8392$, it follows that $\theta^{\rm B}(13,7)>0$. Similarly, $\theta^{\rm B}(24,13)>0$ is the first case where percolation follows for $k=2d-2$, $\theta^{\rm B}(35,19)$ is the first one for $k=2d-3$, and so on.

Although the $k$-UnG model is also an 1-dependent Bernoulli bond percolation model, the same approach unfortunately does not yield any new results in its case. Indeed, here each edge is open with probability $p=k(4d-k)/(4d^2)$. Thus, for $k=d$ it always holds that $p=3/4<0.8457$, while for $k>d$ (and for $d \geq 2$ even for $k \geq 4$) we already know that $\theta^{\rm D}(k,d)=1$. 
\end{remark}

%\subsection{Results for the exclusively unidirectional $k$-neighbor graph}
%Since the $k$-XnG is a subgraph of the $k$-UnG, the following assertion is immediate.
%\begin{prop}
%We have that $\theta^{\rm X}(1,d)=0$ for all $d \geq 1$. 
%\end{prop}



%Using a similar argument, we can also prove the followin
%\begin{lemma}
%We have $\theta(k,d)\leq \theta (k,d')$ whenever $k\geq 1$ and $d\leq d'$.
%\end{lemma}
%\begin{proof}
%The lemma follows from the fact that the restriction of the $k$-nnG in $d'$ dimensions to a $d$-dimensional coordinate plane is stochastically dominated by a $k$-nnG in $d$ dimensions.
%\end{proof}

%\newpage

%%OLD COX-DURRETT BOUND FOR UNDIRECTED/BIDIRECTIONAL, IT HAS ERRORS
%THIS DOESN'T SEEM TO WORK OUT DUE TO PROBLEMS WITH THE MARTINGALES.
%
%With methods analogous to the case of the $k$-UnG, we can additionally prove the following.
%\begin{prop}\label{prop-U}
 %   We have $\theta^{\rm U}(3,4)>0$. Moreover, $\theta^{\rm U}(2,d)>0$ for all $d \geq 5$. \color{red} Need to check issues with lack of independence. \color{black}
%\end{prop}
%\begin{proof}
%We argue analogously to the proof of Theorem~\ref{thm-CoxDurrett}, the only difference being that now the probability of an edge being present in the graph is higher. \color{red} (Add negativity of correlations from the old stuff. \color{black} Namely, we have for an undirected edge $\ell$ with endpoints $x_1,x_2 \in \Z^d$ being nearest neighbors:
%\begin{equation}\label{Uprob}
%\begin{aligned}
%\P&(\ell\text{ is an edge of the $k$-UnG}) = \P((x_1,x_2)\text{ is an edge of the $k$-UnG}) \\ & \quad +\P((x_2,x_1)\text{ is an edge of the $k$-UnG})- \P(\text{both } (x_1,x_2) \text{ and } (x_2,x_1) \text{ are edges of the $k$-UnG} ) \\ &= \frac{k}{2d} + \frac{k}{2d} - \frac{k^2}{4d^2} = \frac{k}{d} \Big( 1 - \frac{k}{4d} \Big),
%\end{aligned}
%\end{equation}
%since the events $\{ (x_1,x_2)\text{ is an edge of the $k$-UnG}\}$ and $\{ (x_2,x_1) \text{ is an edge of the $k$-UnG} \}$ are independent and their probabilities equal $k/(2d)$. Now, using the martingale $W_n=N_n/((\frac{k}{d} ( 1 - \frac{k}{4d} ))^n)$, if $\varrho(d)< \frac{k}{d} ( 1 - \frac{k}{4d} )$, it follows analogously to the proof of Theorem~\ref{thm-CoxDurrett} that $\theta^{\rm U}(k,d)>0$. 
%
%Let us first consider the case $(k,d)=(3,4)$. In the proof of Theorem~\ref{thm-CoxDurrett} we computed that $\varrho(4)=\P(\tau_4<\infty) < 0.5$, and thus
%\[ \varrho(4) < \frac{k}{d} \Big( 1 - \frac{k}{4d} \Big)= \frac{3}{4} \Big( 1- \frac{3}{4 \cdot 4} \Big) = \frac{39}{64} = 0.609375. \] 
%Thus, $\theta^{\rm U}(3,4)>0$.
%
%Let us now focus on the case $(k,d)=(2,5)$. We also computed in the aforementioned proof that $\varrho(5)<0.28$, and thus
%\[ \varrho(5) < \frac{k}{d} \Big( 1 - \frac{k}{4d} \Big) = \frac{2}{5} \Big( 1- \frac{2}{20} \Big) = 0.36. \]
%This implies that $\theta^{\rm U}(2,5)>0$. For $d=6,7,8,9,10$ it in fact holds that the right-hand side of~\eqref{CDUB} is smaller than $\frac{2}{d} (1-2/(4d))$, and for $d \geq 10$ it is again elementary to show that this right-hand side of~\eqref{CDUB} minus $\frac{2}{d} (1-2/(4d))$ stays negative (given that it is negative for $d=9$). \color{red} (Indeed, I think that we can again argue that the rhs of~\eqref{CDUB} for dimension $d+1$ is at most $d/(d+1)$ times the same for dimension $d$, while $\frac{2}{d+1}(1-2/(4d+1))>\frac{d}{d+1} \frac{2}{d}(1-2/(4d))$. \color{black}) Therefore, $\theta^{\rm U}(2,k)>0$ holds for all $k\geq 5$, as asserted.
%\end{proof}
%
%\color{green} If this is fine, the conjecture is proven apart from $(k,d)=(2,2),(2,3),(3,3),(2,4)$. For $(2,4)$ it is not impossible that a Cox--Durrett style bound can work, but we would need $\P(\tau_4=k)$ at least for all $k\leq 8$. So the problem is similar as in the directed case for $k=3$ and $d=4$. \color{black}
%
%
%\begin{prop}\label{prop-B}
%For $d \geq 5$ and $k \geq \sqrt{6 d}$ and for $d \geq 4$ and $k\geq \sqrt{8d}$ we have $\theta^{\rm B}(k,d)>0$.
%\end{prop}
%This implies e.g.\ that $\theta^{\rm B}(5,5)>0$ and $\theta^{\rm B}(6,4)>0$. \color{green} So far, $k=5$ is the smallest $k$ for which we get percolation for some $d$. \color{black}
%\begin{proof}[Proof of Proposition~\ref{prop-B}]
%We again adapt the proof of Theorem~\ref{thm-CoxDurrett}, similarly to how we did it in the proof of Proposition~\ref{prop-U}. \color{red} Also here we have negative correlations -- need to copy it from old stuff. \color{black} For an undirected edge $\ell$ with endpoints $x_1,x_2 \in \Z^d$ being nearest neighbors we have (cf.~\eqref{Uprob})
%\[ \P(\ell \text{ is an edge of the $k$-BnG}) = \P(\text{both } (x_1,x_2) \text{ and } (x_2,x_1) \text{ are edges of the $k$-UnG} )=\frac{k^2}{4d^2}. \]
%Now, we have seen in the aforementioned proofs that $\varrho(d)<\frac{3}{2d}$ for $d \geq 5$ and $\varrho(d)<\frac{2}{d}$ for $d \geq 4$. Consequently, using the martingale $W_n=(4d^2/k^2)^n N_n$, we obtain percolation for $d\geq 5$ whenever 
%\[ \frac{3}{2d} < \frac{k^2}{4d^2}, \]
%which is equivalent to $k \geq \sqrt{6 d}$. For $d \geq 4$ we obtain percolation whenever
%\[ \frac{2}{d} < \frac{k^2}{4d^2}, \]
%in other words, if $k \geq \sqrt{8d}$, as claimed.
%\end{proof}


\section{Discussion and Open Problems}\label{sec:discussion}
\subsection{Relation between discrete and continuum $k$-nearest-neighbor percolation}\label{sec:motivation}
One motivation for studying our particular model is as follows. Consider the continuum nearest-neighbor percolation model with $k=2$, where the intensity of the underlying Poisson point process is such that the expected number of points in a unit ball equals one, and let $Y_i$ denote the position of the $i$-th nearest neighbor for $i\in\{1,2\}$. H\"aggström and Meester show that $|Y_i|$ converges in probability to one as $d$ tends to infinite, as well as that the conditional distribution of $Y_i$, given $|Y_i|=r$, is uniform on the sphere $\{x\in \R^d: |x|=r\}$~\cite[Lemma $3.2$]{HaggMee96}. Moreover, the volume of the intersection of two spheres with radius $r_1,r_2\in(0.9,1.1)$, whose centers are at least $0.9$ units apart, is negligible compared to the volume of either sphere~\cite[Lemma $3.3$]{HaggMee96}. As a result, for large dimensions the continuum nearest-neighbor graph has connections at distance around $1$, which are established almost independently among different pairs of vertices. As such, an approximation of this graph on $\Z^d$ should yield similar behavior when $d$ is large. 

\subsection{Phase transitions from the perspective of complementary relations}\label{sec:phases}
Let the graphs $G,G_1,G_2$ all be directed (respectively all undirected) with the same vertex set. We say that $G_1$ is the complement of $G_2$ w.r.t.\ $G$ if the edge sets of $G_1$ and $G_2$ are disjoint and their union equals the edge set of $G$. For undirected graphs with $n \in \N$ edges, choosing $G=K_n$ yields the usual definition of graph-theoretic complement. Then, the following assertions hold for $1 \leq k < 2d$.
\begin{itemize}
    \item The complement of the $k$-DnG with respect to $\Z^d$ with all directed nearest-neighbor edges is a $(2d-k)$-DnG. (In particular, the $d$-DnG equals its complement in distribution.)
    \item The complement of the $k$-UnG with respect to $\Z^d$ with all undirected nearest-neighbor edges is a $(2d-k)$-BnG. (Hence, the same is also true with the words ,,UnG'' and ,,BnG'' swapped.) 
    \item $\Z^d$ with all directed nearest-neighbor edges is a $2d$-DnG and with all undirected nearest-neighbor edges is a $2d$-UnG and a $2d$-BnG at the same time. Hence, the above two assertions also hold for $k \in \{ 0, 2d\}$, where the 0-DnG is defined as $\Z^d$ as a directed graph with no edges and the 0-UnG and 0-BnG as $\Z^d$ as an undirected graph without edges.
    %\item For $1 \leq k \leq 2d$, the graph-theoretic complement of the $k$-XnG is the disjoint union of two correlated graphs, their marginal distributions being the one of a $k$-BnG and a $2d-k$-BnG. 
\end{itemize}

In what follows, we simply speak about complements of the graphs w.r.t.\ $\Z^d$, which should always be understood in the corresponding (directed resp.\ undirected) sense. Let now $d \in \N$ and $1 \leq k < 2d$. Borrowing a nomenclature known from continuum percolation \cite[Section 7.3]{MeesterRoy1996}, we say that the $k$-UnG in $d$ dimensions
\begin{itemize}
    \item is an \emph{infinite cascade} if $\theta^{\rm U}(k,d)=0$ and $\theta^{\rm B}(2d-k,d)=0$,
    \item it is \emph{supercritical with subcritical complement} if $\theta^{\rm U}(k,d)>0$ and $\theta^{\rm B}(2d-k,d)=0$, %\footnote{\color{red} András: thinking a bit more about it, for classical critical Bernoulli percolation in high dimensions the right term would be ``supercritical with critically non-percolating complement''. So my nomenclature in these definitions may be a bit misleading, we could write ``percolating with nonpercolating complement'', but I still like the current version better. What do you think? \color{black}}
    \item it is \emph{subcritical with supercritical complement} if $\theta^{\rm U}(k,d)=0$ and $\theta^{\rm B}(2d-k,d)>0$,
    \item and it exhibits a \emph{coexistence of phases} $\theta^{\rm U}(k,d)>0$ and $\theta^{\rm B}(2d-k,d)>0$,
\end{itemize}
The same notions can be defined analogously for the $k$-BnG, swapping the roles of $\rm U$ and $\rm B$ everywhere in the above definitions, and for the $k$-DnG, replacing each $\rm U$ and $\rm B$ by a $\rm D$. In words, an infinite cascade means that neither the graph nor its complement (w.r.t.\ $\Z^d$) has an infinite connected component, supercriticality with subcritical complement means that the graph has one but its complement has none, while a coexistence of phases means that both graphs have one.   

Now, for $k=d$, since the $k$-DnG is equal to its own complement w.r.t.\ $\Z^d$ in distribution, either it is an infinite cascade or it exhibits a coexistence of phases. Further, for $k=d$, the complement of the $k$-BnG w.r.t.\ $\Z^d$ equals the $k$-UnG in distribution, and hence the $k$-BnG cannot be supercritical with subcritical complement, while the $k$-UnG cannot be subcritical with supercritical complement. Finally, if the $k$-DnG is supercritical (resp.\ subcritical) with subcritical (resp.\ supercritical) complement, then $k>d$ (resp.\ $k<d$). In the directed case, the fact that $\theta^{\rm D}(k,d)=1$ for $k>d$ implies that an infinite cascade is only possible for $k=d$.

According to our main results, the following assertions hold.
\begin{itemize}
    \item It is trivial that in one dimension, the $1$-DnG is an infinite cascade, and so are the $1$-UnG and the $1$-BnG. Even if there are nontrivial examples for an infinite cascade, there are only finitely many of them. The only possible examples in the directed case are $k=d=2$ and $k=d=3$. However, Theorem~\ref{theorem-3322}  implies that if the $3$-DnG for $d=3$ is an infinite cascade, then so is the $2$-DnG for $d=2$, and as already anticipated, numerical simulations actually suggest that in this case we have a coexistence of phases, see Figure~\ref{fig:2dng} and the corresponding parts of Figures~\ref{fig:boundprob} and \ref{fig:proportion}. For the undirected/bidirectional graph, the possible cases of infinite cascades (with $k$ referring to the undirected and $2d-k$ to the complementary bidirectional graph) are $d=2$ for $k=1$, $3 \leq d \leq 5$ for $k=2$ and $d=4$ for $k=3$. Nevertheless, the $k$-BnG for $k=3=2d-1$ seems to percolate even for $d=2$ according to simulations (see Figure~\ref{fig:3bng}), so that we actually conjecture that none of these cases corresponds to an infinite cascade.
    \item Supercriticality with subcritical complement occurs in all dimensions, both in the directed and undirected cases. Since the $k$-DnG (or the $k$-UnG) percolates for $k\geq 3$ if $d \geq 5$, in these dimensions it can only be subcritical with supercritical component for $k=1,2$, and we know that for $k=1$ this is actually the case, even for $d \geq 2$. The ``sporadic'' open case $k=3, d=4$ that cannot be an infinite cascade could be another example where the $k$-DnG is subcritical with supercritical complement; otherwise we have a coexistence of phases. Also for $k=3$ and $d=4$, but for the UnG, none of the four possibilities can be excluded. For $d=2,3$, we have proven that the $d$-UnG is supercritical with subcritical complement. In Figures~\ref{fig:boundprob} and \ref{fig:proportion} we see that in this case, the probability that the connected component of the origin reaches the boundary of a box of size $2n+1$ and the proportion of vertices in this box belonging to this component both tend to 1 for the $2$-UnG and to 0 for the complementary $2$-BnG rapidly as $n\to\infty$.  For $k \geq 1$ fixed, we have seen that the same holds for the $(2d-k)$-UnG if $d$ is sufficiently large (and for $k=1$ actually for all $d \geq 2$ and for $k=2$ for all $d \geq 3$). 
    \item A coexistence of phases occurs in many high-dimensional cases both in the directed and the undirected--bidirectional model. In the directed case, for $d \geq 5$, since $2d-3 > d$ for such a dimension, it follows that there is a coexistence of phases for $k \in \{3, \ldots, 2d-3 \}$.  We also observe a coexistence of phases for $k=d=4$, which is the smallest-dimensional known example of a coexistence of phases so far.  Also for fixed $k \geq 3$, there is a coexistence of phases for the $k$-UnG for $d$ sufficiently large, and the same holds in general for $k=k(d)$ such that $k(d) < 2d-2 \sqrt{0.5847}d$. The smallest-dimensional example where we verified a coexistence of phases in the undirected--bidirectional model is the $3$-UnG (and the $13$-BnG) in 8 dimensions. %\basc{Shouldn't this be $4$-UnG and $18$-BnG, according to the remark after Corollary~\ref{cor-Bperc}? Also, we can now decrease the dimension to $10$, so that it becomes a $4$-UnG and $16$-BnG? I may be misunderstanding something, so feel free to ignore.} \color{red} András: Well-spotted, thanks, it should be something with a 3-UnG but I didn't choose the smallest possible dimension. $3$-UnG and $13$-BnG in 8 dimensions, I think this is the best that we have now. \color{black}  
    It would be interesting to know if a coexistence of phases can be observed in any of the models for $d=2,3$; a natural candidate mentioned already is the $2$-DnG in 2 dimensions.
\end{itemize}
In this discussion we have not distinguished between the cases when $o \rightsquigarrow \infty$ happens with probability in $(0,1)$ or with probability 1. The latter is the case in the $k$-DnG and $k$-UnG if (and only if!) $k>d$ (cf.\ Lemma~\ref{lem:sureperc}). This does not exclude a coexistence of phases in the directed case (take e.g.\ $k=3$ and $d \geq 5$). It is an interesting open question whether the same holds for the undirected--bidirectional case. We are not aware of any case where $k<d$ and $\theta^{\rm B}(k,d)>0$, which would be necessary for such an assertion.

In classical Bernoulli bond percolation where edges are open independently with probability $p$, it is a well-known result that there is no percolation at the critical point $p_{\rm c}$ for $d=2$, and the same assertion is also known to hold in all sufficiently high dimensions, see e.g.\ \cite{CST} and the references therein. For $p=1/2$, which equals $p_{\rm c}$ in two dimensions, the complement of the percolation model (w.r.t.\ $\Z^d$) is equal to the model itself in distribution, and hence for $d=2$ it provides a nontrivial example for an infinite cascade. Since $p_{\rm c}(d) \sim (2d)^{-1}$, for sufficiently high dimensions we have a coexistence of phases for $p=1/2$; see~\cite{Grimmett1999} and the references therein. With these classical results in mind, the potential coexistence of phases in a self-complementary (directed) graph in two dimensions just as the 2-DnG comes as a surprise. \color{black}


\begin{figure}[ht]
    \centering
    \includegraphics[width=0.95\textwidth]{Images/Boundary_probability}
    \caption{Estimate of $\theta_n^{\square}(2,2)$, the probability that the connected component of the origin reaches the boundary of a box of size $2n+1$ centered around the origin, from $1000$ samples for each $n\in\{5,15,25,35,45,55,65,75,85,95,105,115,125\}$, for the model $2$-$\square$nG and $d=2$. Here, $\square\in\{$D,U,B$\}$.}
    \label{fig:boundprob}
\end{figure}

\begin{figure}[ht]
    \centering
    \includegraphics[width=0.95\textwidth]{Images/Prop_comp_origin}
    \caption{Estimate of the proportion of the number of vertices in a box of size $2n+1$ centered around the origin that belong to the connected component of the origin. An average of $1000$ samples for each $n\in\{5,15,25,35,45,55,65,75,85,95,105,115,125\}$, for the model $2$-$\square$nG and $d=2$. Here, $\square\in\{$D,U,B$\}$.}
    \label{fig:proportion}
\end{figure}


\subsection{Open problems}

Although we managed to prove the occurence or absence of percolation in a wide range of cases for the $k$-$\square$nG model ($\square\in\{\mathrm{D,U,B}\}$), there are also many cases where the techniques used did not provide any conclusive results. We summarise a few cases here that we deem to be interesting.

\begin{itemize}
    \item Do the $2$-DnG and $2$-UnG models percolate in all dimensions? That is, do we have $\theta^{\rm D}(2,d)>0$ and $\theta^{\rm U}(2,d)>0$ for all $d\in \N$? We conjecture this to be the case, see also Figures~\ref{fig:2dng}, \ref{fig:boundprob}, and~\ref{fig:proportion} for numerical calculations and simulations that support this conjecture in two dimensions.
    \item Does $\theta^{\rm D}(3,d)>0$ hold for $d=3,4$, and does $\theta^{\rm U}(3,4)>0$ hold? Together with the previous conjecture and the results proved here, this would imply that $\theta^{\rm U}(k,d)>0$ for all $k,d\geq 2$.
    \item Does Conjecture~\ref{conj:dimensioncouple} hold? This would aid in proving many of the open cases, e.g.\ proving $\theta^{\rm D}(2,2)>0$ would directly imply $\theta^{\rm D}(2,d)>0$ for any $d\geq 2$.
    \item What is the smallest $k=k(d)$ such that $\theta^{\rm B}(k(d),d)>0$ holds for all $d$ sufficiently large? And does the $3$-BnG model percolate in $2$ dimensions? See Figure~\ref{fig:3bng} for a simulation that supports this conjecture.
    \item Is it true that $\theta^{\rm B}(k,d)=0$ whenever $k\leq d$?
\end{itemize}

\subsection{In-percolation and strong percolation}\label{sec-instrong}
As we already mentioned in the introduction, our notion of directed percolation~\eqref{thetaD} corresponds to out-percolation in~\cite{BalisterBollobas13}. In this paper, we did not develop any specific proof techniques for strong or in-percolation. Nevertheless, it is clear that for fixed $d$ and $k$ percolation in the $k$-BnG implies strong percolation, while strong percolation implies both in- and out-percolation in the $k$-DnG, moreover, in- or out-percolation in the $k$-DnG implies percolation in the $k$-UnG, analogously to the continuum case (where these implications were mentioned in~\cite{BalisterBollobas13}). Hence, our positive percolation results for the $k$-BnG yield ones for strong percolation, and our negative (out-)percolation results for the $k$-DnG imply the absence of strong percolation. We have seen that for $d$ large, out-percolation in the $k$-DnG occurs already for $k=3$, while for the $k$-BnG percolation definitely requires $k=\Omega(\sqrt d)$ and perhaps even $k>d$. It is an interesting open question whether strong percolation is closer to directed (out-)percolation than to bidirectional percolation with this respect.

We conjecture that for any $k$ and $d$, in-percolation occurs if and only if out-percolation does, but in this paper we did not intend to verify such a general assertion. It should possibly follow from some general mass-transport type argument, but it is not known either in the continuum case. Without such a result, it seems that in-percolation is in general more difficult to deal with than out-percolation, due to increased combinatorial complexity. E.g., showing that the 1-DnG does not percolate was very easy (cf.\ the proof of Lemma~\ref{lem:surenoperc}), but proving the lack of in-percolation in the same graph (cf.\ the proof of Proposition~\ref{prop-U1}) was already a bit trickier. Further, in the two-dimensional Poisson case~\cite[Proof of Theorem 2]{BalisterBollobas13}, the authors showed directly that 
out-percolation occurs for $k=13$, but for in-percolation for the same $k$, the same method did not work. They derived that their arguments for out-percolation actually imply strong and therefore also in-percolation.

\appendix
\section{Proof of Lemma~\ref{lemma-largedmon}}\label{sec:largedmon}
In this appendix, we carry out the proof of Lemma~\ref{lemma-largedmon}.
\begin{proof}[Proof of Lemma~\ref{lemma-largedmon}]
Recall that~\eqref{largedmon}, i.e.,
\[ R(k,d):=\frac{1}{d} + \frac{1}{d^3} -\frac{1}{d^4} + \sum_{k=3}^d d^{-k} k! +  (2\pi d)^{1/2} \Big( \frac{\e^{-1/13}}{\sqrt{2\pi}} \Big)^d \zeta\Big( \frac{d-1}{2} \Big) - \frac{k}{2d} < 0  \]
holds for any $k\geq 3$ if $d=7$. 
Our goal is to prove the statement of Lemma~\ref{lemma-largedmon}, i.e., that the same holds for any $d \geq 7$ and $k\geq 3$. We check numerically that it also holds for $8 \leq d \leq 11$ and $k=3$ (and thus also for $k\geq 4$ for the same values of $d$): indeed, we have $R(8,3)=-0.0292277$, $R(9,3)=-0.0350912$, $R(10,3)=-0.0367514$ and $R(11,3)=-0.0364418$. 

Now, it is easy to show that for $d \geq 11$, we have
\begin{equation}\label{Rfactor} R(k,d+1) \leq \frac{d}{d+1} R(k,d) + (d+1)^{-(d+1)} (d+1)! \end{equation}
Indeed, it holds that
\[ 
\begin{aligned}
   & R(k,d+1)-(d+1)^{-(d+1)}(d+1)! \\ &=\frac{1}{d+1}- \frac{k}{2(d+1)} + \frac{1}{(d+1)^3} -\frac{1}{(d+1)^4} + \sum_{k=3}^{d} (d+1)^{-k} k! +  (2\pi (d+1))^{1/2} \Big( \frac{\e^{-1/13}}{\sqrt{2\pi}} \Big)^{d+1} \zeta\Big( \frac{d}{2} \Big)
\\    &=\frac{d}{d+1} \Big( \frac{1}{d}- \frac{k}{2d} \Big) + \frac{1}{(d+1)^3} -\frac{1}{(d+1)^4} + \sum_{k=3}^{d} (d+1)^{-k} k! +  (2\pi (d+1))^{1/2} \Big( \frac{\e^{-1/13}}{\sqrt{2\pi}} \Big)^{d+1} \zeta\Big( \frac{d}{2} \Big)
\\ & \leq \frac{d}{d+1} \Big( \frac{1}{d}- \frac{k}{2d} \Big) + \frac{d}{d+1} \bigg[ \frac{1}{d^3} -\frac{1}{d^4} + \sum_{k=3}^d d^{-k} k! +  (2\pi d)^{1/2} \Big( \frac{\e^{-1/13}}{\sqrt{2\pi}} \Big)^d \zeta\Big( \frac{d-1}{2} \Big) \bigg] \\ &=\frac{d}{d+1} R(k,d),
\end{aligned}  \]
where in the last step we used that $\zeta$ is monotone decreasing in $(1,\infty)$ and that for $d \geq 11$,
\begin{equation*} \sqrt{\frac{(d+1)}{ d}} \frac{\e^{-1/13}}{\sqrt{2\pi}} \leq \sqrt{\frac{12}{11}} \frac{\e^{-1/13}}{\sqrt{2\pi}}  \approx 0.385831, \end{equation*}
while $d/(d+1)>11/12$, which is strictly larger. 

Next, let us show that
\begin{equation}\label{newtermbound} (d+1)^{-(d+1)} (d+1)! \leq \frac{d}{d+1} \big( \frac{1}{d^3}-\frac{1}{d^4} \big) - \Big( \frac{1}{(d+1)^3} - \frac{1}{(d+1)^4} \Big) \end{equation}
holds for $d \geq 11$. Note that for such $d$, we have
\[ (d+1)^{-(d+1)} (d+1)! \leq \frac{12!}{(d+1)^{12}} \leq \frac{12!}{12^8} \frac{1}{(d+1)^4} \approx \frac{1.114}{(d+1)^4}, \]
and
\[ \frac{d}{d+1} \Big( \frac{1}{d^3}-\frac{1}{d^4} \Big) - \Big( \frac{1}{(d+1)^3} - \frac{1}{(d+1)^4} \Big) = \frac{(d-1)(d+1)^3}{d^3(d+1)^4} = \frac{2d^3-2d-1}{d^3(d+1)^4} > \frac{1.99}{(d+1)^4} \]
Hence,~\eqref{newtermbound} holds for $d \geq 11$, and thus such $d$ and for any $k\geq 3$,
\[ R(k,d+1) \leq \frac{d}{d+1} R(k,d). \]
Since $R(k,11) < 0$, the lemma now follows by induction over $d$.
\end{proof}


\section*{Acknowledgements}
The authors thank D.~Mitsche, G.~Pete, and B.~Ráth for interesting discussions and comments. BJ and JK acknowledge the financial support of the Leibniz Association within the Leibniz Junior Research Group on \textit{Probabilistic Methods for Dynamic Communication Networks} as part of the Leibniz Competition. BL is supported by the grant GrHyDy ANR-20-CE40-0002. AT  was partially supported by the ERC Synergy under
Grant No.\ 810115 - DYNASNET. 


\printbibliography



\end{document}

